\section{Proof of \autoref{thm: iv mcopi with 1/g scaling}}
\label{app: proof of stochastic theorem}

%%%%%%%%%%%%%%%%%%%%%%%%%%%%%%%%%%%%%%%%%%%%%%%%%%%%%%
%%%%%%%%%%%%%%%%%%%%%%%%%%%%%%%%%%%%%%%%%%%%%%%%%%%%%%
% \subsection{Difference equation}

% \autoref{thm: iv mcopi with 1/g scaling} is essentially an application of the following result.

% \begin{theorem}
% \label{thm: iv mcopi converges as}
% Assume the sampling distributions $\{\sigma_k\}_{k \ge 0}$ are uniform over the actions in each state:
% \begin{align}
% \label{eq: uniform sampling}
% % \sigma_k(s,a) = \sigma_k(s,a')
% \sigma_k(s,a) = \frac{f_k(s)}{|\aset(s)|},
% \qquad 
% \forall \, k \ge 0, \
% \forall \, s \in \sset, \
% \forall \, a \in \aset(s) .
% \end{align}
% % for every time $k \ge 0$, every state $s \in \mathcal S$, and all actions $a, a' \in \mathcal A(s)$.
% Under \autoref{asp: standing}, solutions of IV MCES \eqref{eq: update iv mcopi} converge almost surely to $q_{\pi_*}$.
% \end{theorem}

% This appendix proves \autoref{thm: iv mcopi converges as} 
% and then shows that the same proof applies to \autoref{thm: iv mcopi with 1/g scaling}



Define the perturbation $v_k := \tilde q_k - q_{\pi_k}$ where $\tilde q_k$ is the episode return defined in \eqref{eq: episode return},
as well as the indicator $u_k(s, a) := \mathbf{1}_{\{(s,a)=(s_k,a_k)\}}$ that captures the randomness of asynchronous updates.
Combine these two sources of noise into a single stochastic perturbation $w_k$ defined as
\begin{align}
\label{eq: noise definition appendix}
w_k(s,a) 
&:=
\frac{1}{g_k(a | s)}
\Big( u_k(s,a) v_k(s,a)
+
\big( u_k(s,a) - \sigma_k(s,a) \big) \big( q_{\pi_k}(s,a) - q_k(s,a) \big) \Big) .
\end{align}
Equation \eqref{eq: update iv mcopi} with learning rates as in \eqref{eq: normalised learning rates} is then equivalent, for every state-action pair, to
\begin{align}
% \label{eq: intermediate stochastic step}
\nonumber
q_{k+1}(s,a)
&=
q_k(s,a) + \alpha_k \frac{u_k(s,a)}{g_k(a | s)} \big( q_{\pi_k}(s,a) + v_k(s,a) - q_k(s,a) \big)
\\
\label{eq: stochastic rep of normalised}
    &=
    q_k(s,a) + \alpha_k \big( f_k(s) ( q_{\pi_k}(s,a) - q_k(s,a) ) + w_k(s,a) \big) 
\end{align}
where $\pi_k$ is a deterministic policy that is greedy with respect to $q_k$.
% The sequence $\{\pi_k\}_{k \ge 0}$ thus remains in the finite set of deterministic policies, denoted $\mathcal P$.

\autoref{thm: iv mcopi with 1/g scaling} is an adaptation of Proposition~20 in \cite{Oliviers2026ConvergenceUpdates}. The event construction and policy-improvement argument are unchanged: 
stochastic perturbations cannot consistently undermine the policy improvement driven by the mean-field dynamics,
so $\pi_k$ eventually stabilises at $\pi_*$ and $q_k \to q_{\pi_*}$.


The only modification is the normalized learning rate $\alpha_k/g_k(a\mid s)$ in place of $\alpha_k$. Since $g_k(a\mid s) \in (0,1]$, this rescaling is positive and at least as large as $\alpha_k$, so the induction in Lemma~15 is unaffected. Lemmas~17 and~18 require only that $w_k$ is a martingale difference sequence with bounded second moment; this is verified in the lemma below.



% While several standard frameworks can derive the convergence of a stochastic difference inclusion from the corresponding mean-field dynamics, they each face structural limitations
% when applied to MC-O-PI.

% As hinted in the proof sketch, the convergence argument consists of showing that the stochastic perturbations $\{w_k\}_{k \ge 0}$ cannot consistently prevent the mean-field policy improvement.
% As a result, the policy sequence eventually reaches the optimal policy $\pi_*$ and solutions converge to its action values $q_{\pi_*}$.
% So we start by deriving fundamental properties of these perturbations.
% Unless stated otherwise, operations and relations between functions are elementwise. For instance, given functions $f, g : \xset \to \yset$ we write
% \begin{align}
% &(f + g)(x) = f(x) + g(x), \qquad \forall \, x \in \xset ,
%     \\
% &f \le g \iff f(x) \le g(x), \qquad \forall \, x \in \xset .
% \end{align}
% Furthermore, every norm is computed as $\|f\| := \sup_{x \in \xset} |f(x)|$.
% % The stochastic perturbations $w_k$ sat

% %%%%%%%%%%%%%%%%%%%%%%%%%%%%%%%%%%%%%%%%%%%%%%%%%%%%%%
% %%%%%%%%%%%%%%%%%%%%%%%%%%%%%%%%%%%%%%%%%%%%%%%%%%%%%%
% % \subsection{Basic properties}

\begin{lemma}
% [Martingale difference noise]
\label{thm: properties of noise}
% Under initial- or first-visit updates, 
There exists a constant $c_w \in [0, \infty)$ 
such that
\begin{align}
\label{eq: condition on combined noise for proof}
\mathbb E[w_k(s,a) \mid \mathcal F_k] = 0,
\ \ \
\mathbb E[w_k^2(s,a) \mid \mathcal F_k] \le c_w(1+\|q_k\|^2),
\ \ \
\forall \, k \ge 0, \,
\forall \, s \in \mathcal S, \,
\forall \, a \in \mathcal A(s) .
\end{align}
% for every time $k \ge 0$, every state $s \in \sset$ and every action $a \in \aset(s)$,
% \begin{align}
% \label{eq: muw_mds}
% &\E\!\left[w_k(s,a) \mid \mathcal F_k\right] = 0,
% \\
% \label{eq: muw_second_moment}
% &\E\!\left[w_k^2(s,a) \mid \mathcal F_k\right]
% \le c_w(1+\|q_k\|^2) .
% \end{align}
\end{lemma}
\begin{proof}
Fix a time $k\ge 0$, a state $s \in \mathcal S$ and an action $a \in \mathcal A(s)$.
Let $\mathcal F_k'$ represent the available information right after selecting the policy $\pi_k$ but before sampling the initial state-action pair $(s_k, a_k)$, so $\mathcal F_k \subseteq \mathcal F_k'$.

Under initial-visit updates, $\tilde q_k$ is an unbiased estimate of $q_{\pi_k}$ with uniformly bounded variance
\cite{Singh1996ReinforcementTraces}.
Thus there exists a constant $c_v \in [0,\infty)$ such that
\begin{align}
\label{eq: properties episode noise}
\mathbb E\!\left[v_k(s,a) \mid \mathcal F_k \right] = 0 ,
\qquad
\mathbb E\!\left[v_k^2(s,a) \mid \mathcal F_k \right] \le c_v,
\qquad
\forall \, k \ge 0 .
\end{align}
Further, 
% $\mathbb E[u_k(s,a) \mid \mathcal F_k] = \sigma_k(s,a) = f_k(s) g_k(a | s)$ under initial-visit updates, 
% and 
$u_k(s,a) v_k(s,a) = 0$ on $\{u_k(s,a)=0\}$.
Hence,
\begin{align}
\label{eq: noise unbiased part 1}
\mathbb E\big[ u_k(s,a) v_k(s,a) \mid \mathcal F'_k \big]
=
\mathbb E \Big[ 
    \mathbf{1}_{\{u_k(s,a)=1\}} \ 
    \mathbb E \big[ v_k(s,a) \mid \mathcal F'_k, u_k(s,a) = 1 \big]
    \mid \mathcal F'_k \Big]
    = 0 .
\end{align}
Also,
$\E[u_k(s,a) \mid \mathcal F'_k] = \sigma_k(s,a)$ 
and
% given that $k$, $s$ and $a$ are fixed,
$\sigma_k(s,a)$ and $(q_{\pi_k}(s,a)-q_k(s,a))$ are $\mathcal F'_k$-measurable. It follows that
\begin{multline}
\mathbb E \!\left[ 
\big( u_k(s,a) - \sigma_k(s,a) \big) \big( q_{\pi_k}(s,a) - q_k(s,a) \big)
\mid \mathcal F'_k
\right]
\\
\label{eq: noise unbiased part 2}
=
\big(
\mathbb E [ u_k(s,a) \mid \mathcal F'_k ]
- \sigma_k(s,a)
\big)
\big( q_{\pi_k}(s,a)-q_k(s,a) \big)
=
0 .
\end{multline}
Since $g_k(a | s)$ is $\mathcal F'_k$-measurable,
dividing \eqref{eq: noise unbiased part 1} and \eqref{eq: noise unbiased part 2} by $g_k(a | s)$ and summing the results yields $\mathbb E[w_k(s,a) \mid \mathcal F'_k] = 0$.
Applying the tower property then gives
\begin{align}
\label{eq: noise is unbiased}
\mathbb E [ w_k(s,a) \mid \mathcal F_k ]
= \mathbb E \big[ \mathbb E [ w_k(s,a) \mid \mathcal F'_k] \mid \mathcal F_k \big]
= 0 .
\end{align}

Finally,
$g_k^2(a | s) \ge \sigma_{\min}^2$ by the strict exploring starts condition in \eqref{eq: strict exploring starts},
% $|\mathcal A(s)| \ge \min_{s} |\mathcal A(s)| =: |\mathcal A_{\min}|$,
$u_k^2(s,a) \le 1$ 
and $(u_k(s,a) - \sigma_k(s,a))^2 \le 1$ for every state-action pair.
Using $(x+y)^2 \le 2 x^2 + 2 y^2$ on \eqref{eq: noise definition appendix} thus yields
\begin{align*} 
w_k^2(s,a) 
&\le
\frac{1}{\sigma_{\min}^2} \big( 2 v_k^2(s,a) + 2 (q_{\pi_k}(s,a)-q_k(s,a) )^2 \big)
\\
&\le
\frac{1}{\sigma_{\min}^2} \big( 2 v_k^2(s,a) + 4 q_{\pi_k}^2(s,a) + 4 q_k^2(s,a) \big) .
\end{align*}
Since all action-values are bounded,
there exists a constant $c_\polset \in [0,\infty)$ such that $\|q_\pi\|^2 \leq c_\polset$ for every policy $\pi$. 
As a result, using \eqref{eq: properties episode noise} yields
\begin{align}
\nonumber
\mathbb E \!\big[ w_k^2(s,a) \mid \mathcal F_k \big]
% &\le
% \frac{1}{\sigma_{\min}^2 |\mathcal A_{\min}|^2}  \left( 2 c_v + 2 \, \mathbb E \!\left[ ( q_{\pi_k}(s,a)-q_k(s,a) )^2 \mid \mathcal F_k \right] \right)
% \\
% % \label{eq: noise has bounded variance}
% \nonumber
&\le
\frac{2}{\sigma_{\min}^2} ( c_v + 2 c_\polset + 2 \|q_k\|^2 ) .
\\
\label{eq: noise has bounded variance}
&\le
\max\left\{ \frac{2}{\sigma_{\min}^2} ( c_v + 2 c_\polset), \frac{4}{\sigma_{\min}^2} \right\} ( 1 + \|q_k\|^2 ) .
\end{align}
Equations \eqref{eq: noise is unbiased} and \eqref{eq: noise has bounded variance} confirm that the perturbations $\{w_k\}_{k \ge 0}$ satisfy the conditions in \eqref{eq: condition on combined noise for proof}.
\end{proof}

% %%%%%%%%%%%%%%%%%%%%%%%%%%%%%%%%%%%%%%%%%%%%%%%%%%%%%%
% %%%%%%%%%%%%%%%%%%%%%%%%%%%%%%%%%%%%%%%%%%%%%%%%%%%%%%
% % \subsection{Bounded solutions}

% Given these properties of $w_k$, 
% we can apply standard martingale and stochastic approximation arguments to establish fundamental properties of \eqref{eq: stochastic rep of normalised}.
% For instance, if the target action values $q_{\pi_k}$ were fixed,
% solutions would converge almost surely to this fixed target by Proposition~4.1 or Example~4.3 in \cite{Bertsekas1996Neuro-DynamicProgramming}.
% As a result, solutions of the switched system are almost surely bounded.


% \begin{lemma}
% % [Bounded limit set]
% \label{thm: bounded limit sets}
% Consider the setting of \autoref{thm: iv mcopi with 1/g scaling}.
% Let $\polset_\infty$ denote the set of policies that appear infinitely often in a sequence $\{\pi_k\}_{k \ge 0}$ generated by \eqref{eq: stochastic rep of normalised}.
% % \begin{align}
% %     \polset_\infty \coloneqq \bigcap_{t = 0}^{\infty} \text{cl} \{\pi_k : k \geq t \} ,
% % \end{align}
% The corresponding sequence $\{q_k\}_{k \ge 0}$ converges almost surely to the set
% \begin{align*}
%     \Big\{
%     q : 
%     \min_{\pi \in \polset_\infty} q_\pi(s,a) \leq q(s,a) \leq \max_{\pi \in \polset_\infty} q_\pi(s,a), \forall s \in \sset, \forall a \in \aset(s)
%     \Big\} .
% \end{align*}
% % Consequently,
% % Let $\{q_k\}_{k\ge 0}$ be a solution of \eqref{eq: iv mcopi in stochastic theorem}.
% % There exists an almost surely finite constant $c_Q \in [0, \infty)$ such that for all $k \geq 0$
% % \begin{align*}
% %     \| q_k \|_\infty^2 < c_Q .
% % \end{align*}
% \end{lemma}

% \begin{proof}
% Let $\{q_k\}_{k \ge 0}$ and $\{\pol_k\}_{k \ge 0}$ be solutions of \eqref{eq: stochastic rep of normalised} with initial condition $q_0$, learning rates $\{\alpha_k\}_{k \ge 0}$, state-sampling distributions $\{f_k\}_{k \ge 0}$ and stochastic perturbations $\{w_k\}_{k \ge 0}$.
    
% By assumption, there exists a finite time $K$ such that
% the sequence $\{\pol_k\}_{k \ge K}$ only contains policies in $\polset_\infty$.
% % the estimate $q_k$ is updated towards action values $q_{\pi_k}$ where $\pi_k \in \polset_\infty$.
% Thus for all $k \ge K$, each state-action pair $(s,a)$ of the estimate $q_k$ is updated towards a value $q_{\pol_k}(s,a)$ that satisfies
% \begin{align}
%     \min_{\pi \in \polset_\infty} q_\pi(s,a) \leq q_{\pi_k}(s,a) \leq \max_{\pi \in \polset_\infty} q_\pi(s,a) .
% \end{align}
% Therefore, define the sequences $\{x_k\}_{k \ge K}$ and $\{y_k\}_{k \ge K}$ as
% \begin{align}
%     x_\kpo &= x_k + \alpha_k 
%     \left(
%     % \sigma_k(s,a) 
%     f_k(s) \bigg( \min_{\pi \in \polset_\infty} q_\pi(s,a) - x_k \bigg) + w_k(s,a) \right) ,
%     \\
%     y_\kpo &= y_k + \alpha_k 
%     \left(
%     % \sigma_k(s,a) 
%     f_k(s) \bigg( \max_{\pi \in \polset_\infty} q_\pi(s,a) - y_k \bigg) + w_k(s,a) \right) ,
% \end{align}
% and initialised at $x_K = y_K = q_K(s,a)$. From time $K$ onwards we have
% \begin{align*}
%     x_k \leq q_k(s,a) \leq y_k .
% \end{align*}
% Standard results show that $x_k$ and $y_k$ converge almost surely to $\min_{\pi \in \polset_\infty} q_\pi(s,a)$ and $\max_{\pi \in \polset_\infty} q_\pi(s,a)$, respectively~\citep[Proposition~4.1, Example~4.3]{Bertsekas1996Neuro-DynamicProgramming}.
% % , because the noise satisfies conditions \eqref{eq: muw_mds} and \eqref{eq: muw_second_moment}.
% As a result, we conclude that, almost surely,
% \begin{align*}
%     \liminf_{k \to \infty} q_k(s,a)
%     \ge \min_{\pi \in \polset_\infty} q_\pi(s,a) ,
%     % \leq q_k(s,a) 
%     \quad
%     \limsup_{k \to \infty} q_k(s,a)
%     \le \max_{\pi \in \polset_\infty} q_\pi(s,a) ,
%     \quad
%     \forall \, s \in \sset , \
%     \forall \, a \in \aset(s) .
% \end{align*}
% % for every state-action pair.
% \end{proof}

% %%%%%%%%%%%%%%%%%%%%%%%%%%%%%%%%%%%%%%%%%%%%%%%%%%%%%%
% %%%%%%%%%%%%%%%%%%%%%%%%%%%%%%%%%%%%%%%%%%%%%%%%%%%%%%
% % \subsection{Policy transitions}

% Furthermore, given solutions would converge almost surely to $q_{\pi_k}$ if these values were fixed in \eqref{eq: stochastic rep of normalised},
% the asymptotic behaviour of \eqref{eq: stochastic rep of normalised} hinges on how these action values evolve over time.
% To isolate these changes, we introduce \emph{transition times}.
% We adopt the convention $\inf \varnothing = \infty$.

% \begin{definition}
% \label{def: transition times}
% Given a sequence of policies $\{\pi_k\}_{k \ge 0}$,
% define the transition times $\{t_n\}_{n \ge 0}$ as
% \begin{align*}
% t_0 \coloneqq 0 ,
% \qquad
% t_{n+1} \coloneqq \inf\{k>t_n : q_{\pi_k} \neq q_{\pi_{t_n}} \} .
% \end{align*}
% % with the convention that $\inf \varnothing = \infty$.
% Thus, if $t_n = \infty$, then $t_{n+i} = \infty$ for all $i \ge 1$ .
% These times decompose the sequence $\{\pi_k\}_{k \ge 0}$ into blocks of constant action values, such that $q_{\pi_k} = q_{\pi_{t_n}}$ for all $k \in \{t_n, \dots, t_{n+1}-1 \}$.
% \end{definition}

% % Clearly, if $\pi_{t_n}$ is suboptimal, then $t_{n+1}<\infty$.
% % Otherwise the target would remain equal to \(q_{\pi_{t_n}}\) forever, forcing convergence to \(q_{\pi_{t_n}}\).
% % However a suboptimal policy is not greedy with respect to its own action values.
% % Thus no solution can get arbitrarily close to $q_{\pi_{t_n}}$ without causing the greedy policy to improve. Therefore no suboptimal block can persist indefinitely.

% \begin{lemma}
% \label{thm: suboptimal blocks finite}
% If $\pi_{t_n}$ is not optimal, then $t_{n+1}<\infty$ almost surely.
% \end{lemma}
% \begin{proof}
% Let $\pi := \pi_{t_n}$.
% Suppose for contradiction that $t_{n+1}=\infty$. The target in \eqref{eq: stochastic rep of normalised} is then fixed at $q_\pi$ forever, and $q_k \to q_\pi$ almost surely by Proposition~4.1 in \cite{Bertsekas1996Neuro-DynamicProgramming}.
% However, since $\pi$ is suboptimal, it is not greedy with respect to $q_\pi$. Thus there exists a state-action pair $(s,a)$ such that $q_\pi(s,a) > q_\pi(s,\pi(s))$.
% As a result, as $q_k$ approaches $q_\pi$,
% every greedy policy $\pi_k$ with respect to $q_k$ satisfies
% $q_\pi(s,\pi_k(s)) > q_\pi(s,\pi(s))$
% for all sufficiently large $k$.
% The Policy Improvement Theorem then indicates that $q_{\pi_k} \ne q_\pi$, contradicting $t_{n+1}=\infty$.
% \end{proof}

% % 
% % Consequently, if there exists a finite index $n$ such that $t_{n} = \infty$, the associated solution converges to the action values of the last bounded transition time, which must be optimal.
% % Conversely, any suboptimal asymptotic behaviour requires an infinite sequence of bounded transition times.


% We rely on the following strict version of the of the Policy Improvement Theorem at transition times.
% \begin{lemma}
% \label{thm: strict improvement theorem}
% Given policies \(\pi,\pi'\in\mathcal P\), suppose
% \begin{align}
% \label{eq: PIT condition}
% q_\pi(s,\pi'(s)) \ge q_\pi(s,\pi(s))
% \qquad \forall s\in\mathcal S .
% \end{align}
% Then \(v_{\pi'} \ge v_\pi\). Moreover, if \(q_{\pi'}\neq q_\pi\), then the improvement is strict:
% there exists a state \(s\in\mathcal S\) where
% \[
% v_{\pi'}(s)>v_\pi(s).
% \]
% \end{lemma}
% \begin{proof}
% By the Policy Improvement Theorem, inequality \eqref{eq: PIT condition} implies that $v_{\pi'} \ge v_\pi$.
% Suppose that $v_{\pi'} = v_{\pi}$. Then for every state-action pair $(s,a)$,
% \begin{align*}
% q_{\pi'}(s,a)
% = r(s,a) + \gamma \sum_{s' \in \mathcal S} p (s' \mid s,a) v_{\pi'}(s')
% = r(s,a) + \gamma \sum_{s' \in \mathcal S} p (s' \mid s,a) v_{\pi}(s')
% = q_{\pi}(s,a) .
% \end{align*}
% As a result, if $q_{\pi'} \neq q_\pi$, there exists at least one state $s$ where $v_{\pi'}(s) > v_{\pi}(s)$, implying that policy $\pi'$ is better than $\pi$.
% \end{proof}


% %%%%%%%%%%%%%%%%%%%%%%%%%%%%%%%%%%%%%%%%%%%%%%%%%%%%%%
% %%%%%%%%%%%%%%%%%%%%%%%%%%%%%%%%%%%%%%%%%%%%%%%%%%%%%%
% \subsection{Lower bound on the probability of policy improvement at transitions}

% Consider the sequences $\{q_k\}_{k\ge 0}$ and $\{\pi_k\}_{k\ge 0}$ generated by \mcopi, and let $\{t_n\}_{n \ge 0}$ be the associated transition times.
% Over each interval $\{t_n, \dots, t_{n+1}\}$, we compare the evolution of $q_{k}$ at $(s,\pi_{t_n}(s))$ and pairs $(s,a)$ that would worsen the policy $\pi_{t_n}$, namely
% \[
% \mathcal B_{\pi_{t_n}} := \big\{ (s,a) : s \in \mathcal S, \ a \in \mathcal A(s), \ q_{\pi_{t_n}}(s,a) < q_{\pi_{t_n}}(s, \pi_{t_n}(s)) \big\} .
% \]
% % We build an event that guarantees that, 
% We show that, with a fixed probability and for sufficiently large $n$,
% for every $k \in \{t_n, \dots, t_{n+1}\}$
% the greedy policy satisfies $(s, \pi_k(s)) \notin \mathcal B_{\pi_{t_n}}$.
% % either $(s, \pi_{t_{n+1}}(s)) \notin \mathcal B_{\pi_{t_n}}$ for every state $s$,
% % or $t_{n+1} = \infty$.
% % The first case implies that
% % \begin{align}
% % \label{eq: goal of proof stoch improvement}
% % q_{\pi_{t_n}} ( s, \pi_{t_{n+1}}(s) )
% % \ge
% % q_{\pi_{t_n}} ( s,\pi_{t_n}(s) )
% % \end{align}
% % and therefore that the policy $\pi_{t_{n+1}}$ is better than $\pi_{t_n}$ by \autoref{thm: strict improvement theorem}.
% % The second case means that, by \autoref{thm: suboptimal blocks finite}, $\pi_{t_n} = \pi_*$ and thus $q_k$ will converge to the associated optimal action values.
% % Overall, either the policy improves, or it is already optimal and stays optimal.
% As a result, if $t_{n+1}<\infty$ we obtain
% \begin{align}
% \label{eq: goal of proof stoch improvement}
% q_{\pi_{t_n}} ( s, \pi_{t_{n+1}}(s) )
% \ge
% q_{\pi_{t_n}} ( s,\pi_{t_n}(s) )
% \end{align}
% and therefore policy $\pi_{t_{n+1}}$ is better than $\pi_{t_n}$ by \autoref{thm: strict improvement theorem}.
% Or if $t_{n+1} = \infty$,
% \autoref{thm: suboptimal blocks finite} implies that $\pi_{t_n} = \pi_*$ and $q_k$ converges to the associated optimal action values.
% In other words, the policy either improves, or it is already optimal and stays optimal.


% % By \autoref{def: transition times}, the target action values are fixed at $q_{\pi_{t_n}}$ for all times $ k \in \{t_n,\dots,t_{n+1}-1\}$.
% % We know that if $\pi_{t_n}$ is suboptimal then $t_{n+1}$ is almost surely finite, and hence $\pi_{t_{n+1}}$ is well defined.
% % The following result establishes a uniform lower bound on the probability that policy $\pi_{t_{n+1}}$ is better than $\pi_{t_{n}}$.

% % Technically, discounted or episodic MDPs can have several optimal policies.
% % We denote the set of all such optimal policies by $\mathcal P_*$. By definition, all optimal policies share the same value functions, denoted $v_{\best{\pi}}$ and $q_{\best{\pi}}$.
% % Furthermore, at least one policy in $\mathcal P_*$ is deterministic.

% % The main point to proving \autoref{thm: fixed prob that bad actions dont become greedy} is to construct a nonzero-probability event that ensures that


% \begin{proposition}
% \label{thm: fixed prob that bad actions dont become greedy}
%     Let $\{\pi_k\}_{k\ge 0}$ be a sequence of policies generated by
%     system \eqref{eq: stochastic rep of normalised}
%     under the setting of \autoref{thm: iv mcopi with 1/g scaling}.
%     Let $\{t_n\}_{n\ge 0}$ be the associated transition times.
%     There exist a constant $p>0$ and an almost surely finite time $K$ such that
%     for every $n \in \mathbb Z_{\ge 0}$ with $K \le t_n < \infty$,
%     \begin{align}
%     \label{eq: porb better action value}
%     \mathbb P \big(
%     \forall \, s \in \sset, \
%     \forall \, k \in \{t_n, \dots , t_{n+1} \}
%     \ : \
%     (s, \pi_k(s)) \notin \mathcal B_{\pi_{t_n}}
%     \mid \mathcal F_{t_n}
%     \big) \ge p .
%     \end{align}
% \end{proposition}


% We track the distance between $(s,\pi_{t_n}(s))$ and pairs in $\mathcal B_{\pi_{t_n}}$ using gap variables.

% \begin{definition}[Gap variables]
% \label{def: gap variables}
% % Let $\{t_n\}_{n\ge 0}$ be the transition times associated with the policy sequence generated by system \eqref{eq: stochastic rep} under the setting of \autoref{thm: iv mcopi converges as}.
% Fix a finite transition index $n \in \mathbb Z_{\ge 0}$
% and write $\pi := \pi_{t_n}$.
% % For each state $s \in \mathcal S$ and action $a \in \mathcal A(s)$, the gap is initialised at time $t_n$ as
% For each $k \in \{t_n, \dots, t_{n+1}\}$, define
% \begin{align*}
% % \label{eq: def gap initialisation}
% % d_{t_n}(s,a) = q_{t_n}(s,\pi(s)) - q_{t_n}(s,a) ,
% d_{k}(s,a) = q_{k}(s,\pi(s)) - q_{k}(s,a) .
% \end{align*}
% While $k < t_{n+1}$, each gap evolves as
% \begin{align}
% \label{eq: def gap evolution}
% d_{k+1}(s,a) = d_k(s,a) + \alpha_k \left( f_k(s)(d_\pi(s,a) - d_k(s,a)) + \xi_k(s,a) \right)
% \end{align}
% % where we used
% % \begin{align*}
% %     \mu_k(s) \coloneqq 
% %     % \sigma_k(s,\pi(s)) = \sigma_k(s,a) = 
% %     \frac{f_k(s)}{|\aset(s)|},
% % \end{align*}
% with the target gap
% \begin{align*}
% % \label{eq: def policy gap}
%     d_\pi(s,a) \coloneqq q_{\pi}(s,\pi(s)) - q_{\pi}(s,a) ,
% \end{align*}
% and the gap noise
% \begin{align*}
% % \label{eq:xi-def}
%     \xi_k(s,a) &\coloneqq
% w_{k}(s,\pi(s)) - w_{k}(s,a) .
% % \\
% % \nonumber
% % &\: = u_k(s,\pi(s))v_k(s,\pi(s)) - u_k(s,a) v_k(s,a)
% % \\
% % \nonumber
% % &\: \quad + \bigl( u_k(s,\pi(s)) - \mu_k(s) \bigr) \bigl( q_\pi(s,\pi(s)) - q_k(s,\pi(s)) \bigr)
% % \\
% % \nonumber
% % &\: \quad - \bigl( u_k(s,a) - \mu_k(s) \bigr) \bigl( q_\pi(s,a) - q_k(s,a) \bigr) .
% \end{align*}
% \end{definition}

% % Note that we are not saying that the policy does not change between $t_n$ and $t_{n+1}-1$. Instead we just say that it does not matter because the target action values remain fixed and, as long as we can guarantee that no action from $\mathcal B_{\pi_{t_n}}$ is sampled during that time frame that is enough, i.e.
% % \[
% % q_{k}(s,\pi_k(s)) \ge q_{k}(s,\pi_{t_n}(s)) \ge q_{k}(s,a)
% % \]

% % Even though $w_k(s,\pi(s))$ and $w_k(s,a)$ are not independent, the gap noise $\xi_k(s,a)$ is unbiased and its conditional variance is eventually bounded by a constant.

% \begin{lemma}
% \label{thm: properties noise gap}
% There exist a constant $c_\xi \in [0,\infty)$ 
% % such that for every \mcopi{} solution there is 
% and an almost surely finite time $K_\xi$ such that
% for every transition index $n \in \mathbb Z_{\ge 0}$ satisfying $t_n \ge K_\xi$,
% \begin{align*}
% % \label{eq:xi-mean}
% \mathbb E [ \xi_{k}(s,a) \mid \mathcal F_k ] = 0 ,
% % \\
% % \label{eq:xi-var}
% \quad
% \mathbb E [ \xi_{k}^2(s,a) \mid \mathcal F_k ] \le c_\xi ,
% \quad
% \forall \, k \in \{t_n, \dots, t_{n+1}-1\}, \
% \forall \, s \in \mathcal S, \
% \forall \, a \in \mathcal A(s) .
% \end{align*}
% % for every $k \in \{t_n, \dots, t_{n+1}-1\}$, every $s \in \mathcal S$, and every $a \in \mathcal A(s)$.
% % For every time $k \ge 0$,
% % every state $s \in \mathcal S$,
% % and every action $a \in \mathcal A(s)$,
% % \begin{align}
% % \label{eq:xi-mean}
% % &\mathbb E\!\left[ \xi_{k}(s,a) \mid \mathcal F_k \right] = 0 .
% % \end{align}
% % Moreover, there exists a constant $c_\xi>0$ 
% % % such that for every \mcopi{} solution there is 
% % and an almost surely finite time $K_\xi$ such that
% % for every time $k \ge K_\xi$,
% % every state $s \in \mathcal S$,
% % and every action $a \in \mathcal A(s)$,
% % \begin{align}
% % \label{eq:xi-var}
% % &\mathbb E\!\left[ \xi_{k}^2(s,a) \mid \mathcal F_k \right] \le c_\xi .
% % \end{align}
% \end{lemma}
% \begin{proof}
% \autoref{thm: properties of noise}
% % requires that $w_k(s,a)$ is unbiased.
% directly yields
% \[
% \E\!\big[ \xi_{k}(s,a) \mid \mathcal F_k \big] 
% =
% \E\!\big[ w_{k}(s,\pi_{t_n}(s)) \mid \mathcal F_k \big] - \E\!\big[ w_{k}(s,a) \mid \mathcal F_k \big] 
% = 
% 0 .
% \]

% Further,
% \autoref{thm: bounded limit sets} implies that
% % $q_k$ converge almost surely to the set $\bigl\{ q \in \mathcal Q : q_{\worst{\pi}} \le q \le q_{\best{\pi}} \bigr\}$, where $\overline \pi$ and $\pi_*$ are the worst and best policies of the MDP, respectively.
% % Thus, 
% there exists an almost surely finite time $K_{\xi}$ such that $\| q_k \| \le c_{\mathcal P} + 1$ for every $k \ge K_{\xi}$,
% with $c_{\mathcal P} := \max_{\pi}\| q_\pi \|_\infty$.
% Substituting this into \eqref{eq: condition on combined noise for proof} yields
% \[
% \mathbb E[w_k^2(s,a) \mid \mathcal F_k] 
% \le c_w(1+\|q_k\|^2)
% \le c_w(1+(c_{\mathcal P} + 1)^2)
% \qquad
% \forall \, k \ge K_{\xi}, \
% \forall \, s \in \mathcal S, \
% \forall \, a \in \mathcal A(s).
% \]
% Using $(x-y)^2 \le 2 ( x^2 + y^2)$ then shows that,
% % proves
% % \eqref{eq:xi-var}
% for every $k \ge K_{\xi}$, every $s \in \mathcal S$, and all $a, a' \in \mathcal A(s)$,
% \begin{align*}   
% \mathbb E\!\left[ \bigl(w_k(s,a) - w_k(s,a') \bigr)^2 \ \big| \ \mathcal F_k \right]
% &\le
% 2 \left( \mathbb E\!\left[ w_k^2(s,a) \mid \mathcal F_k \right] 
% + \mathbb E\!\left[ w_k^2(s,a') \mid \mathcal F_k \right] \right)
% \\
% &\le
% 4 \, c_w \left( 1 + (c_{\mathcal P}+1)^2 \right)
% =: c_\xi .
% \end{align*}
% % for every $k \ge K$, every $s \in \mathcal S$, and all $a, a' \in \mathcal A(s)$.
% % Hence, \eqref{eq:xi-var} holds with $c_\xi := 4 \, c_w \left( 1 + (c_{\mathcal P}+1)^2 \right)$ and $K_\xi = K$.
% \end{proof}


% % We show that, since each gap is non-negative at initialisation because $\pi_{t_n}$ is greedy with respect to $q_{t_n}$,
% % and (2) the target $d_{\pi_{t_n}}$ is strictly positive for each pair in $\mathcal B_{\pi_{t_n}}$,
% % there is a fixed probability that no pair in $\mathcal B_{\pi_{t_n}}$ is greedy at time $t_{n+1}$, and hence that we obtain \eqref{eq: goal of proof stoch improvement}.

% Over each interval $\{t_n, \dots, t_{n+1}\}$,
% we construct three events that guarantee that the gap $d_k$ of each pair in $\mathcal B_{\pi_{t_n}}$ is strictly positive up to time $t_{n+1}$,
% and consequently that no such pair is greedy.
% Specifically, we fix a margin $\delta>0$.
% The \textit{seed} event $\mathcal E^s_{t_n}$ first forces each gap at least up to $O(\alpha_k)$.
% The \textit{growth} event $\mathcal E^g_{t_n}$ then drives each gap from $O(\alpha_k)$ to $\delta$.
% Finally, the \textit{lock-in} event $\mathcal E^l_{t_n}$ ensures that after reaching $\delta$, each gap remains strictly positive up to time $t_{n+1}$.


% Concretely, define the minimum target gap across all policies
% \[
% d_{\min}
% \coloneqq
% \min_{\pi\in\mathcal P}\ \inf_{(s,a)\in \mathcal B_\pi} d_\pi(s,a) > 0 .
% \]
% % with the convention that $\min\varnothing=\infty$.
% Using the constants $\sigma_{\min}$ and $\rho_\alpha$ from
% \eqref{eq: strict exploring starts}--\eqref{eq: comparability of learning rates},
% we choose $\delta$ such that
% \begin{equation}
% \label{eq:delta-choice}
% 0<\delta<\min\left\{\frac{d_{\min}}{8},\,\frac{\rho_\alpha \sigma_{\min} d_{\min}}{32}\right\}.
% \end{equation}
% This allows us to define the constants
% \begin{align}  
% \label{eq:c-def}
% % \label{eq:cs-def}
% c_s \coloneqq \left( \frac{d_{\min}}{4} - \delta \right) > 0,
% % \\
% % \label{eq:cg-def}
% \quad
% c_g \coloneqq \sigma_{\min} ( d_{\min} - 2 \delta) > 0,
% % \\
% % \label{eq:c-def}
% \quad
% c \coloneqq \min \left\{ c_s, c_g \right\} > 0.
% \end{align}


% %%%%%%%%%%%%%%%%%%%%%%%%%%%%%%%%%%%%%%%%%%%%%%%%%%%%%%%%%%%%%%%
% \paragraph{Seed event pushes gap up to $\bm{O(\alpha_k)}$.}

% Let $M \in \mathbb N$ denote a finite time horizon whose value will be fixed later.
% For any time $k \in \{t_n, \ldots, t_{n+1}-1\}$, we define the seed event differently if the value of the greedy action is under- or overestimated.
% % Concretely, for any index $n \in \mathbb Z_{\ge 0}$,
% % any time $k \in \{t_n, \dots, t_{n+1}-1\}$,
% % any state $s \in \mathcal S$,
% % and any action $a \in \mathcal A(s)$,
% Namely, if $q_{\pi_{t_n}}(s,\pi_{t_n}(s)) - q_{k}(s,{\pi_{t_n}}(s)) \ge d_{\min}/2$,
% we update the pair $(s,{\pi_{t_n}}(s))$ for $M$ consecutive steps and require that the noise at each step is not too negative. We gather these conditions in the event
% \begin{align*}
%     \mathcal E^+_k(s,a)
%     \coloneqq 
%     \bigcap_{k'=k}^{k+M-1}
%     \left\{
%     u_{k'}(s,\pi_{t_n}(s)) = 1, \
%     v_{k'}(s,\pi_{t_n}(s)) \ge -\frac{d_{\min}}{4}
%     \right\} .
% \end{align*}
% On the other hand, if $q_{\pi_{t_n}}(s,\pi_{t_n}(s)) - q_{k}(s,{\pi_{t_n}}(s)) < d_{\min}/2$,
% we update $(s,a)$ and require that the noise is not too positive:
% \begin{align*}
%     \mathcal E^-_k(s,a)
%     \coloneqq 
%     \bigcap_{k'=k}^{k+M-1}
%     \left\{
%     u_{k'}(s,a) = 1, \
%     v_{k'}(s,a) \le \frac{d_{\min}}{4}
%     \right\} .
% \end{align*}
% Altogether, we obtain the seed event
% \begin{multline}
% \label{eq: def good prefix}
%     \mathcal E^s_k(s,a) \coloneqq 
%     \left( \left\{ q_{\pi_{t_n}}(s,\pi_{t_n}(s)) - q_{k}(s,{\pi_{t_n}}(s)) \ge \frac{d_{\min}}{2} \right\} \cap \mathcal E^+_k(s,a) \right) 
%     \\
%     \bigcup 
%     \left( \left\{ q_{\pi_{t_n}}(s,\pi_{t_n}(s)) - q_{k}(s,\pi_{t_n}(s)) < \frac{d_{\min}}{2} \right\} \cap \mathcal E^-_k(s,a) \right) .
% \end{multline}
% We also define the hitting time
% \[
% \tau^s_k(s,a) \coloneqq \inf \big\{ k' \in \{k+1, \dots, k+M \} : q_{k'}(s,\pi_{t_n}(s))-q_{k'}(s,a) \ge \delta \} ,
% \]
% and the accumulated learning rate mass between times $k$ and $k'>k$ as
% \[
% A(k,k') \coloneqq \sum_{i=k}^{k'-1} \alpha_{i} .
% \]
% % with the convention that $A(k, k) = 0$.
% We now show that, on $\mathcal E^s_k(s,a)$, ``bad'' gaps grow at least up to $O(\alpha_k)$ in $M$ steps.

% \begin{lemma}[Seed event]
% \label{thm: event to reach a_k margin}
% Consider the setting of \autoref{thm: fixed prob that bad actions dont become greedy}.
% For every index $n \in \mathbb Z_{\ge 0}$,
% every pair $(s,a) \in \mathcal B_{\pi_{t_n}}$
% and every time $k \in \{t_n, \dots, t_{n+1}-1\}$,
% the joint event
% \begin{align}
% \label{eq: joint seed event}
% \mathcal E^s_k(s,a)
% \cap
% % \big\{0 \leq q_{k}(s,\pi_{t_n}(s)) - q_{k}(s,a) < \delta \big\}
% \big\{0 \leq d_k(s,a) < \delta \big\}
% \end{align}
% guarantees that 
% % for every $k'$ satisfying 
% % % $k < k' \le \big( \tau^s_k(s,a) \wedge (k+M) \wedge t_{n+1} \big)$,
% % $k < k' \le \min\{ \tau^s_k(s,a), k+M, t_{n+1} \}$,
% \begin{align}
% \label{eq: lower bound during good prefix}
%     d_{k'}(s,a)
%     % q_{k'}(s,\pi_{t_n}(s)) - q_{k'}(s,a)
%     > c_s A(k,k') > 0
%     \qquad 
%     \forall \, k < k' \le \min\{ \tau^s_k(s,a), k+M, t_{n+1} \} .
% \end{align}
% % for every $k'$ that satisfies $k < k' \le \big( \tau^s_k(s,a) \wedge (k+M) \wedge t_{n+1} \big)$.
% Moreover, 
% % the probability of $\mathcal E^s_k(s,a)$ is lower bounded by
% \begin{align}
% \label{eq: probability reach a_k margin}
%     \mathbb P ( \mathcal E^s_k(s,a) \mid \mathcal F_k, \pi_{t_n} ) \ge \left( \sigma_{\min} \frac{d_{\min}^2}{d_{\min}^2 + 16 c_v} \right)^M > 0 .
% \end{align}
% \end{lemma}

% \begin{proof}
% Fix a transition index $n \in \mathbb Z_{\ge 0}$ and let $\pi := \pi_{t_n}$.
% Fix a pair $(s,a) \in \mathcal B_\pi$,
% define the sequence $\big\{ d_k(s,a) \big\}_{k=t_n}^{t_{n+1}}$ as in \autoref{def: gap variables}, and write $d_k := d_k(s,a)$ and $d_\pi := d_\pi(s,a)$ for brevity.
% Fix a time $k \in \{t_n, \dots, t_{n+1}-1\}$ and write $\tau^s_k := \tau^s_k(s,a)$.

% % We split the proof along the two mutually exclusive cases that make up $\mathcal E^s_k(s,a)$.


% When $q_\pi(s,\pi(s)) - q_{k}(s,\pi(s)) \ge d_{\min}/2$,
% % On 
% event $\mathcal E^+_k(s,a)$
% % only $(s, \pi(s))$ is updated.
% % during $M$ consecutive steps,
% % and $v_{k'}(s,\pi(s)) \ge -d_{\min}/4$ for every $k' \in \{k, \dots, k+M-1\}$.
% implies that for every $k'$ satisfying
% % $k \le k' < \big((k+M) \wedge t_{n+1} \big)$,
% $k \le k' < \min\{ k+M, t_{n+1} \}$,
% \begin{align}
% \nonumber
%     d_{k'+1}
%     &= d_{k'} + 
%     % \alpha_{k'} 
%     \frac{\alpha_{k'}}{g_{k'}(\pi(s) | s)}
%     \big( q_\pi(s,\pi(s)) - q_{k'}(s,\pi(s)) + v_{k'}(s,\pi(s)) \big)
%     \\
% \label{eq: plug noise into update of seed stage case 1}
%     &\ge d_{k'} + 
%     % \alpha_{k'} 
%     \frac{\alpha_{k'}}{g_{k'}(\pi(s) | s)}
%     \left( q_\pi(s,\pi(s)) - q_{k'}(s,\pi(s)) - \frac{d_{\min}}{4} \right) .
% \end{align}
% Given that $0 \le d_{k} < \delta$ and $d_{k'} < \delta$ while $k'< \tau^s_k$,
% it follows that 
% \begin{align}
% % \nonumber
% \label{eq: lower bound update step in seed stage case 1}
%     q_\pi(s,\pi(s)) - q_{k'}(s,\pi(s))
%     = q_\pi(s,\pi(s)) - d_{k'} + d_{k} - q_{k}(s,\pi(s))
%     % \\
%     > \frac{d_{\min}}{2} - \delta > 0
% \end{align}
% as long as
% % $k \le k' < \big( \tau^s_k \wedge (k+M) \wedge t_{n+1} \big)$,
% $k \le k' < \min\{ \tau^s_k , k+M, t_{n+1} \}$.
% Equations \eqref{eq: plug noise into update of seed stage case 1} and \eqref{eq: lower bound update step in seed stage case 1}
% thus indicate that, on \eqref{eq: joint seed event},
% % for every $k'$ satisfying
% % $k \le k' < \min\{ \tau^s_k , k+M, t_{n+1} \}$,
% % the gap strictly increases at each time $k'$ in the restricted range:
% \begin{align}
% \label{eq: final lower bound case 1}
%     d_{k'+1}
%     > d_{k'} +
%     \alpha_{k'}
%     \left( \frac{d_{\min}}{4} - \delta \right)
%     = d_{k'} + \alpha_{k'} c_s
%     % > 0 .
%     \qquad
%     \forall \, k \le k' < \min\{ \tau^s_k , k+M, t_{n+1} \} .
% \end{align}

% Similarly, when $q_\pi(s,\pi(s)) - q_{k}(s,\pi(s)) < d_{\min}/2$
% % event $\mathcal E^-_k(s,a)$ guarantees that
% % only $(s, a)$ is updated during $M$ consecutive steps, and $v_{k'}(s,a) \le d_{\min}/4$ for every $k' \in \{k, \dots, k+M-1\}$.
% we obtain
% \begin{align}
% % \nonumber
%     d_{k'+1}
%     % &= d_{k'} + \alpha_{k'}(q_{k'}(s,a) - q_\pi(s,a) - v_{k'}(s,a))
%     % \\
% \label{eq: plug noise into update of seed stage case 2}
%     &\ge d_{k'} + 
%     % \alpha_{k'} 
%     \frac{\alpha_{k'}}{g_{k'}(a | s)}
%     \left( q_{k'}(s,a) - q_\pi(s,a) - \frac{d_{\min}}{4} \right)
% \end{align}
% as long as
% % $k \le k' < \big((k+M) \wedge t_{n+1} \big)$,
% $k \le k' < \min\{ k+M, t_{n+1}\}$.
% Since $q_{k}(s,\pi(s))$ is unchanged on $\mathcal E^-_k(s,a)$
% and $d_{k'} < \delta$ while $k' < \tau^s_k$,
% it follows that
% \begin{align}
% % \nonumber
%     q_{k'}(s,a) - q_\pi(s,a)
%     = d_\pi - \big( q_\pi(s,\pi(s)) - q_{k}(s,\pi(s)) \big) - d_{k'}
%     % \\
% \label{eq: lower bound update step in seed stage case 2}
%     % > d_{\min} - \frac{d_{\min}}{2} - \delta
%     % \ge \frac{d_{\min}}{2} - \delta 
%     > \frac{d_{\min}}{2} - \delta
%     > 0
% \end{align}
% as long as
% % $k \le k' < \big( \tau^s_k \wedge (k+M) \wedge t_{n+1} \big)$,
% $k \le k' < \min\{ \tau^s_k, k+M, t_{n+1} \}$.
% Hence \eqref{eq: plug noise into update of seed stage case 2} and \eqref{eq: lower bound update step in seed stage case 2} confirm that, on \eqref{eq: joint seed event},
% % the gap strictly increases at each time $k'$ in the restricted range:
% \begin{align}
% \label{eq: final lower bound case 2}
%     d_{k'+1}
%     > d_{k'} + 
%     \alpha_{k'}
%     \left( \frac{d_{\min}}{4} - \delta \right)
%     = d_{k'} + \alpha_{k'} c_s
%     % > 0 .
%     \qquad
%     \forall \, k \le k' < \min\{ \tau^s_k , k+M, t_{n+1} \} .
% \end{align}
% By induction on \eqref{eq: final lower bound case 1} or \eqref{eq: final lower bound case 2},
% % and shifting the index by 1 to include time step $k'+1$,
% we obtain
% \begin{align}
% \label{eq: increase in d^b case B}
%     d_{k'}
%     > d_{k} + c_s A(k,k')
%     \ge c_s A(k,k') > 0
%     \qquad
%     \forall \, k < k' \le \min\{ \tau^s_k , k+M, t_{n+1} \} .
% \end{align}
% % for every $k'$ satisfying
% % $k < k' \le \min\{ \tau^s_k , k+M, t_{n+1} \}$,
% % as long as $k < k' \le \big( \tau^s_k \wedge (k+M) \wedge t_{n+1} \big)$, 
% which proves \eqref{eq: lower bound during good prefix}.

% Finally, 
% % we lower-bound the probability of $\mathcal E^s_k(s,a)$.
% given that $v_k$ is unbiased and its variance is bounded by $c_v$ uniformly over the policies,
% Cantelli's inequality implies that for every $k' \ge 0$, every $s \in \mathcal S$ and every $a \in \mathcal A(s)$,
% \begin{align}
% &\mathbb P \! \left( v_{k'}(s,a) \ge -\frac{d_{\min}}{4} \,\middle|\, \mathcal F_{k'}, u_{k'}(s,a)=1 \right) \ge \frac{d_{\min}^2}{d_{\min}^2 + 16 c_v},
% \label{eq:noise-lower}
% \\
% &\mathbb P \! \left( v_{k'}(s,a) \le \frac{d_{\min}}{4} \,\middle|\, \mathcal F_{k'}, u_{k'}(s,a)=1 \right) \ge \frac{d_{\min}^2}{d_{\min}^2 + 16 c_v}.
% \label{eq:noise-upper}
% \end{align}
% Since the case choice in \eqref{eq: def good prefix} is measurable on $\{\mathcal F_{k}, \pi_{t_n}\}$ and the minimum sampling probability of any state-action pair is $\sigma_{\min}$ by \eqref{eq: strict exploring starts},
% % under strict exploring starts,
% repeated conditioning over $M$ steps yields
% \begin{align}
% \label{eq: prob of good prefix}
% \mathbb P \big( \mathcal E^s_k(s,a) \mid \mathcal F_{k} , \pi_{t_n} \big) \ge \left( \sigma_{\min} \frac{d_{\min}^2}{d_{\min}^2 + 16 c_v} \right)^M > 0 .
% \end{align}
% Since $n$, $(s,a)$ and $k$ were chosen arbitrarily from their respective sets, the result holds uniformly.
% \end{proof}

% %%%%%%%%%%%%%%%%%%%%%%%%%%%%%%%%%%%%%%%%%%%%%%%%%%%%%%%%%%%%%%%
% \paragraph{Growth event pushes gap from $\bm{O(\alpha_k)}$ to $\bm{\delta}$}

% % Let $N(k)$ denote the minimum number of steps required to accumulate a particular amount of learning rate mass:
% Define the growth window $N(k)$ as
% \begin{align}
% N(k)
% % \coloneqq
% % \inf \left\{ i \ge 1 : 
% % A(k,k+i)
% % \ge \frac{4 \delta}{c_g} \right\} .
% \:=
% \left\lceil \frac{4 \delta}{c_g \rho_\alpha \alpha_k} \right\rceil . 
% \end{align}

% \begin{lemma}
% \label{thm: growth window}
% % Under condition \eqref{eq: comparability of learning rates},
% There exists a finite time $K_N$ such that for every $k \ge K_N$,
% \[
% N(k) \le \frac{1}{\alpha_{k}},
% \qquad 
% A(k,k + N(k)) \ge \frac{4 \delta}{c_g} .
% \]
% \end{lemma}
% \begin{proof}
% The definitions of $\delta$ and $c_g$ in \eqref{eq:delta-choice}--\eqref{eq:c-def} ensure that
% \begin{multline*}
%     \frac{4 \delta}{c_g \rho_\alpha}
%     < 
%     \frac{4 \rho_\alpha \sigma_{\min} d_{\min}}{32 \sigma_{\min}(d_{\min} - 2 \delta) \rho_\alpha}
%     =
%     \frac{d_{\min}}{8 (d_{\min} - 2 \delta)}
% %    \\
%     <
%     \frac{d_{\min}}{8(d_{\min} - d_{\min}/4)}
%     =
%     \frac{1}{6}
%     < 1 .
% \end{multline*}
% The Robbins--Monro conditions in \eqref{eq: robbins monro conditions} imply that $\alpha_k \to 0$. So there exists a finite $K$ such that $4\delta / c_g \rho_\alpha + \alpha_k < 1$ for all $k \ge K$,
% and hence
% \[
% N(k) < \frac{4 \delta}{c_g \rho_\alpha \alpha_k} + 1 < \frac{1}{\alpha_k} .
% \]
% % Let
% % \[
% % r_k := \left\lceil \frac{4 \delta}{c_g \rho_\alpha \alpha_k} \right\rceil .
% % \]
% % Since $r_k \le 1/ \alpha_k$,
% The comparability property of the learning rates in \eqref{eq: comparability of learning rates} then guarantees that $A(k, k+N(k)) \ge N(k) \rho_\alpha \alpha_{k} \ge 4 \delta / c_g$ for every $k \ge K_n := \max\{K, K_\alpha\}$.
% % \[
% % A(k, k+r_{k}) \ge r_{k} \rho_\alpha \alpha_{k} \ge 4 \delta / c_g .
% % \]
% % As a result, there exists a finite time $K_N \ge \max\{K, K_\alpha\}$ such that $N(k) \le r_k \le 1/\alpha_{k}$ for all $k \ge K_N$.
% \end{proof}

% For any $k \in \{t_n, \dots, t_{n+1}-1\}$,
% % any state $s \in \mathcal S$
% % and any action $a \in \mathcal A(s)$,
% define the accumulated stochastic perturbations between $k$ and $k'>k$ as
% \[
% \Xi(k, k' \mid s,a) 
% \coloneqq 
% \sum_{i=k}^{k'-1} \alpha_{i} \xi_i(s,a)
% =
% \sum_{i=k}^{k'-1} \alpha_{i} \big( w_i(s, \pi_{t_n}(s))-w_i(s,a) \big) ,
% \]
% % with the convention $\Xi(k, k \mid s,a) = 0$.
% define the growth event
% % For any index $n \in \mathbb Z_{\ge 0}$,
% % any time $k \in \{t_n, \dots, t_{n+1}-1\}$,
% % any state $s \in \mathcal S$
% % and any action $a \in \mathcal A(s)$,
% as
% \[
% \mathcal E^g_k(s,a)
% \coloneqq
% \bigcap_{k'=k+1}^{k+N(k)}
% \left\{
% \Xi(k,k' \mid s,a) > - \frac{c_s}{2} \rho_\alpha M \alpha_k - \frac{c_g}{2} A(k,k')
% \right\} ,
% \]
% % for any $M \in \mathbb N$,
% and define the hitting time as
% \[
% \tau^g_k(s,a) \coloneqq \inf \! \big\{ k' \in \{k+1, \dots, k+N(k) \} : q_{k'}(s,\pi_{t_n}(s)) - q_{k'}(s,a) \ge \delta \big\} .
% \]
% We now show that, on $\mathcal E^g_k(s,a)$, 
% perturbations cannot prevent ``bad'' gaps from reaching $\delta$ in $N(k)$ steps.
% % significantly offset the mean-field drift.
% % the gap $d_k(s,a)$ stays strictly positive until $\min\{\tau^g_k(s,a), t_{n+1}\}$, and reaches $\delta$ within $N(k)$ steps, meaning

% \begin{lemma}[Growth event]
% \label{thm: from a_k margin to delta}
% Consider the setting of \autoref{thm: fixed prob that bad actions dont become greedy}
% and let $K_\xi$ and $K_N$ be the finite times obtained in \autoref{thm: properties noise gap} and \autoref{thm: growth window}, respectively.
% For every index $n \in \mathbb Z_{\ge 0}$ such that $t_n \ge K_N$,
% every pair $(s,a) \in \mathcal B_{\pi_{t_n}}$
% and every time $k \in \{t_n, \dots, t_{n+1}-1\}$,
% the joint event
% \begin{align}
% \label{eq: joint growth event}
% \mathcal E^g_k(s,a) \cap \big\{ c_s \rho_\alpha M \alpha_k \leq d_k(s,a) < \delta \big\}
% % \mathcal E^g_k(s,a) \cap \big\{ c_s \rho_\alpha M \alpha_k \leq q_{k}(s,\pi_{t_n}(s)) - q_{k}(s,a) < \delta \big\}
% \end{align}
% guarantees that 
% % the gap $d_k(s,a)$ stays strictly positive until $\min\{\tau^g_k(s,a), t_{n+1}\}$, and reaches $\delta$ within $N(k)$ steps, meaning
% % for every $k'$ satisfying
% % $k \le k' \le \min \{ \tau^g_k(s,a), t_{n+1} \}$,
% \begin{align}
% \label{eq: growth event gap remains positive}
%     d_{k'}(s,a) > 0
%     % q_{k'}(s,\pi_{t_n}(s)) - q_{k'}(s,a) > 0,
%     \qquad \forall \, k \le k' \le \min \{ \tau^g_k(s,a), t_{n+1} \} , 
% \end{align}
% and
% \begin{align}
% \label{eq: hit time event a_k to delta}
%     \min \{ \tau^g_k(s,a) , t_{n+1} \} \le k + N(k) .
% \end{align}
% Moreover, if $t_n \ge K_\xi$,
% \begin{align}
% \label{eq: prob from a_k reach delta}
%     \mathbb P (\mathcal E^g_k(s,a) \mid \mathcal F_k, \pi_{t_n})
%     \ge 1 - \frac{8 c_\xi}{c^2 \rho_\alpha^2} \sum_{j = 0}^{\infty} \frac{2^j}{(M + 2^j)^2} .
% \end{align}
% \end{lemma}
% \begin{proof}
% Fix an $n \in \mathbb Z_{\ge 0}$ such that $t_n \ge K_N$,
% and let $\pi := \pi_{t_n}$.
% Fix a pair $(s,a) \in \mathcal B_\pi$,
% define the sequence $\big\{ d_k(s,a) \big\}_{k=t_n}^{t_{n+1}}$ as in \autoref{def: gap variables}, and write $d_k := d_k(s,a)$, $d_\pi := d_\pi(s,a)$ and $\xi_k := \xi_k(s,a)$ for brevity.
% Fix a time $k \in \{t_n, \dots, t_{n+1}-1\}$ and write $\tau^g_k := \tau^g_k(s,a)$.

% For every $k' \in \{k, \dots, \tau^g_k\}$ we have $d_{k'} \le \delta < d_{\min} \le d_\pi$, and hence
% \[
% f_{k'}(s) (d_\pi - d_{k'}) 
% \ge 
% \sigma_{\min} (d_{\min} - \delta)
% >
% \sigma_{\min} (d_{\min} - 2 \delta)
% = c_g .
% \]
% Thus summing \eqref{eq: def gap evolution} from $k$ to any $k'$ satisfying $k < k' \le \min\{\tau^g_k, t_{n+1}\}$ yields
% \[
% d_{k'}
% >
% d_{k} + c_g A(k,k') + \Xi(k,k' \mid s,a) .
% \]
% % $k' \in \{k, \dots, t_{n+1}-1\}$
% % the gap satisfies
% % \[
% % d_{k'+1}
% % =
% % d_{k'} + \alpha_{k'} \big( \mu_{k'}(s) (d_\pi - d_{k'}) + \xi_{k'} \big) ,
% % \]
% % where the definition of $\delta$ guarantees that $d_\pi \ge d_{\min} > 2 \delta$,
% % and the strict exploring starts ensure that $\mu_{k'}(s) \ge \sigma_{\min}$. 
% % Since $d_{k'} < \delta$ before the hitting time $\tau^g_k$,
% % we can bound the update for all $k \le k' < ( \tau^g_k \wedge t_{n+1} )$ as
% % \begin{align*}
% % d_{k'+1}
% % >
% % d_{k'} + \alpha_{k'} \left( \sigma_{\min}(d_{\min} - \delta) + \xi_{k'} \right)
% % % \\
% % % &>
% % % d_{k'} + \alpha_{k'} \left( \sigma_{\min}(d_{\min} - 2 \delta) + \xi_{k'} \right)
% % % \\
% % >
% % d_{k'} + \alpha_{k'} \left( c_g + \xi_{k'} \right) ,
% % \end{align*}
% % using $c_g \coloneqq \sigma_{\min} ( d_{\min} - 2 \delta)$.
% % Summing from $k$ to $k'$ then yields
% % \[
% % d_{k'}
% % >
% % d_{k} + c_g A(k,k') + \Xi(k,k' \mid s,a) .
% % \]
% As a result, the joint event in \eqref{eq: joint growth event} implies that as long as $k < k' \le \min\{ \tau^g_k , t_{n+1}, k+N(k) \}$,
% \begin{align}
% d_{k'}
% &>
% c_s \rho_\alpha M \alpha_k + c_g A(k,k') - \frac{c_s}{2} \rho_\alpha M \alpha_k - \frac{c_g}{2} A(k,k')
% \nonumber
% \\
% \label{eq: gap remains positive up to tau'_k}
% &=
% \frac{c_s}{2} \rho_\alpha M \alpha_k + \frac{c_g}{2} A(k,k') > 0 .
% \end{align}
% % which proves \eqref{eq: growth event gap remains positive}.
% % 
% Now assume that $t_{n+1} > k+N(k)$, otherwise $\min\{ \tau^g_k, t_{n+1} \} \le k+N(k)$ is immediate.
% Suppose for contradiction that $\tau^g_k = \infty$ because $d_k$ does not reach $\delta$ within $N(k)$ steps. 
% Since $k \ge t_n \ge K_N$, \autoref{thm: growth window} and equation \eqref{eq: gap remains positive up to tau'_k} yield
% \[
% d_{k+N(k)}
% >
% \frac{c_s}{2} \rho_\alpha M \alpha_k + \frac{c_g}{2} A(k,k+N(k)) 
% >
% % \frac{c_g}{2} A(k,k+N(k)) 
% % \ge
% \frac{c_g}{2} \frac{4 \delta}{c_g} 
% > 
% \delta .
% \]
% This contradicts $\tau^g_k = \infty$. Hence $\tau^g_k \le k+N(k)$, which proves \eqref{eq: hit time event a_k to delta}; and with \eqref{eq: gap remains positive up to tau'_k} also gives \eqref{eq: growth event gap remains positive}.

% % This contradiction implies that $\tau^g_k \le k+N(k)$, which proves \eqref{eq: hit time event a_k to delta} and, together with \eqref{eq: gap remains positive up to tau'_k}, proves \eqref{eq: growth event gap remains positive}.

% % Then \eqref{eq: gap remains positive up to tau'_k} yields
% % \[
% % d_{k+N(k)}
% % >
% % \frac{c_s}{2} \rho_\alpha M \alpha_k + \frac{c_g}{2} A(k,k+N(k)) 
% % >
% % \frac{c_g}{2} A(k,k+N(k)) .
% % \]
% % However,

% % since the definition of $N(k)$ guarantees that $A(k,k+N(k)) \ge 4 \delta/c_g$, we obtain a contradiction: $d_{k+N(k)}
% % >
% % \frac{c_s}{2} \rho_\alpha M \alpha_k + 2 \delta
% % > \delta .$
% % % \[
% % % d_{k+N(k)}
% % % >
% % % \frac{c_s}{2} \rho_\alpha M \alpha_k + 2 \delta
% % % > \delta .
% % % \]
% % Thus, $\tau^g_k \le k+N(k)$, which proves \eqref{eq: hit time event a_k to delta}.

% Finally,
% % we lower-bound the probability of $\mathcal E^g_k(s,a)$.
% the complement event $\overline{\mathcal E^g_k(s,a)}$ implies that for some $i \in \{1, \dots , N(k) \}$,
% \begin{align*}
% \bigl| \Xi(k,k+i \mid s,a) \bigr|
% \ge \frac{c_s}{2} \rho_\alpha M \alpha_k + \frac{c_g}{2} A(k,k+i)
% % \]
% % Using $c \coloneqq \min\{c_s,c_g\}$ and the comparability property (since $i \le N(k) \le 1/\alpha_k$ for $k \ge t_n \ge K_N$), this simplifies to
% % \[
% % \bigl| \Xi(k,k+i \mid s,a) \bigr|
% \ge \frac{c}{2} \big( \rho_\alpha M \alpha_k + A(k,k+i) \big) .
% % \ge \frac{c}{2} \rho_\alpha \alpha_k ( M + i ).
% \end{align*}
% Choosing $j \in \mathbb N$ such that $2^j \le i \le 2^{j+1}$ 
% and using the comparability of the learning rates \eqref{eq: comparability of learning rates}
% yields the bound
% \begin{align}
%     \max_{1 \le \ell \le \min \{ 2^{j+1} , N(k) \}}
%     |\Xi(k,k + \ell \mid s,a)|
%     \ge \frac{c}{2} \big( \rho_\alpha M \alpha_k + A(k,k + 2^j) \big)
%     \ge \frac{c}{2} \rho_\alpha \alpha_k (M + 2^j) .
% \end{align}
% So if $t_n \ge K_\xi$,
% % By \autoref{thm: properties noise gap}, there exists an almost surely finite $K_\xi$ such that 
% the sequence $\big\{\Xi(k, k+\ell \mid s,a)\big\}_{\ell \ge 1}$ is a square-integrable martingale adapted to the filtration $\{\mathcal F_{k+\ell}\}_{\ell \ge 1}$. So, by Doob's martingale inequality,
% \begin{align}
% \label{eq: prob of not E^g_k}
%     \mathbb P
%     \left( 
%     \overline{\mathcal E^g_k(s,a)} \mid \mathcal F_{k}, \pi_{t_n}
%     \right)
%     % \\
%     % \nonumber
%     \le
%     \sum_{j = 0}^{\infty} \frac{\mathbb E \left[ \big( \Xi(k, \min\{ k+2^{j+1}, k+N(k) \} \mid s,a) \big)^2 \mid \mathcal F_{k}, \pi_{t_n} \right]}{\left( \dfrac{c}{2} \rho_\alpha \alpha_k ( M + 2^j) \right)^2} .
% \end{align}
% \autoref{thm: properties noise gap} and the non-increasing sequence $\{\alpha_k\}_{k \ge 0}$ imply that
% \begin{align*}
% \mathbb E \left[ \big( \Xi(k, \min\{ k+2^{j+1}, k+N(k)\} \mid s,a) \big)^2 \mid \mathcal F_{k}, \pi_{t_n} \right]
% \le
% c_\xi \sum_{\ell=0}^{2^{j+1}-1} \alpha^2_{k+\ell}
% \le
% c_\xi 2^{j+1} \alpha^2_k .
% \end{align*}
% Substituting this back into \eqref{eq: prob of not E^g_k} proves \eqref{eq: prob from a_k reach delta} because
% \begin{align}
%     \mathbb P
%     \left( 
%     \overline{\mathcal E^g_k(s,a)} \mid \mathcal F_{k}, \pi_{t_n}
%     \right)
%     \le
%     \frac{8 c_\xi}{c^2 \rho_\alpha^2} \sum_{j = 0}^{\infty} \frac{2^j}{(M+ 2^j)^2}.
% \end{align}
% Since $n$, $(s,a)$ and $k$ were chosen arbitrarily from their respective sets, the result holds uniformly.
% \end{proof}

% % \autoref{thm: from a_k margin to delta} confirms that on the event $\mathcal E^g_k \cap \{c \rho_\alpha M \alpha_k \leq q_k(s,\pi_{t_n}(s))-q_k(s,a) < \delta \}$, the gap safely hits $\delta$ in at most $N(k)$ steps while remaining strictly positive throughout the ascent.

% %%%%%%%%%%%%%%%%%%%%%%%%%%%%%%%%%%%%%%%%%%%%%%%%%%%%%%%%
% \paragraph{Lock-in event keeps gap positive after reaching $\bm{\delta}$}

% Define the lock-in event as
% \begin{align}
% \label{eq: definition lock in event single gap}
%     \mathcal E^l_{t_n}(s,a) \coloneqq \left\{ \sup_{k \ge t_n} \bigl| \Xi(t_n, k \mid s,a) \bigr| < \frac{\delta}{8} \right\} .
% \end{align}
% We now show that, on $\mathcal E^l_{t_n}(s,a)$,
% once a ``bad'' gap reaches $\delta$, 
% it remains strictly positive indefinitely.
% % perturbations can never drive it below zero anymore.
% % For any $n \in \mathbb Z_{\ge 0}$,
% % any state $s \in \mathcal S$
% % and any action $a \in \mathcal A(s)$,


% \begin{lemma}[Lock-in event]
% \label{thm: event noise remains small}
% Consider the setting of \autoref{thm: fixed prob that bad actions dont become greedy} and let $K_\xi$ be the finite time obtained in \autoref{thm: properties noise gap}.
% For every index $n \in \mathbb Z_{\ge 0}$,
% every pair $(s,a) \in \mathcal B_{\pi_{t_n}}$
% and every time $k \in \{t_n, \dots, t_{n+1}-1\}$,
% the joint event
% \[
% \mathcal E^l_{t_n}(s,a) \cap 
% % \big\{ q_{k}(s,\pi_{t_n}(s)) - q_{k}(s,a) \ge \delta \big\}
% \big\{ d_k(s,a) \ge \delta \big\}
% \]
% guarantees that
% % the gap remains strictly positive until the next transition time, meaning
% % for every $k' \in \{k, \dots, t_{n+1}\}$,
% \begin{align}
% \label{eq: gap remains strictly positive}
% % q_{k'}(s,\pi_{t_n}(s)) > q_{k'}(s,a) 
% d_{k'}(s,a) > 0
% \qquad 
% \forall \, k' \in \{k, \dots, t_{n+1}\} .
% \end{align}
% Moreover, if $t_n \ge K_\xi$,
% % there exists an almost surely finite time $K_\xi$ such that for every $t_n \ge K_\xi$,
% \begin{equation}
% \label{eq:probability-small-remaining-noise}
% \mathbb P\!\left(\mathcal E^l_{t_n}(s,a)\mid \mathcal F_{t_n}, \pi_{t_n} \right)
% \ge
% 1-\frac{64 c_\xi}{\delta^2}\sum_{i=t_n}^\infty \alpha_i^2 .
% \end{equation}
% \end{lemma}
% \begin{proof}
% Fix an $n \in \mathbb Z_{\ge 0}$
% and let $\pi := \pi_{t_n}$.
% Fix a pair $(s,a) \in \mathcal B_\pi$,
% define the sequence $\big\{ d_k(s,a) \big\}_{k=t_n}^{t_{n+1}}$ as in \autoref{def: gap variables}, and write $d_k := d_k(s,a)$, $d_\pi := d_\pi(s,a)$ and $\xi_k := \xi_k(s,a)$ for brevity.
% Fix a time $k \in \{t_n, \dots, t_{n+1}-1\}$.

% % For every time $k \in \{t_n, \dots, t_{n+1}-1\}$, the gap is updated as
% % \begin{align*}
% % d_{k+1}
% % =
% % \big( 1 - \alpha_{k} \mu_{k}(s) \big) d_{k} + \alpha_{k} \mu_{k}(s) d_\pi + \alpha_{k} \xi_{k} .
% % \end{align*}
% % To unroll this recursion, define the contraction factor between steps $k$ and $k'>k$ as
% % \[
% % \Gamma(k,k') := \prod_{i=k}^{k'-1} \big( 1 - \alpha_{i} \mu_{i}(s) \big) \in(0,1].
% % \]
% % with the convention that $\Gamma(k,k) = 1$.


% Define the contraction factor between times $k$ and $k'>k$ as
% \[
% \Gamma(k,k') := \prod_{i=k}^{k'-1} \big( 1 - \alpha_{i} f_{i}(s) \big) \in(0,1] 
% \]
% with the convention that $\Gamma(k,k) = 1$.
% Unrolling \eqref{eq: def gap evolution} from $k$ to any $k' \in \{k,\dots,t_{n+1}\}$ yields
% \begin{align}
% \nonumber
% d_{k'}
% &=
% \Gamma(k,k') d_{k}
% +
% \bigl(1-\Gamma(k,k') \bigr) d_\pi
% + \sum_{i=k}^{k'-1} \Gamma(i+1,k') \alpha_i \xi_i
% \\
% \label{eq: induction for abel proof}  
% &\ge
% \delta + \sum_{i=k}^{k'-1} \Gamma(i+1,k') \alpha_i \xi_i
% \end{align}
% because $\Gamma(k,k') \in(0,1]$,
% $d_\pi \ge d_{\min} > \delta$
% % $\Gamma(k,k') \in (0,1]$,
% and, by assumption, $d_{k} \ge \delta$.
% % we have
% % \begin{align}
% % \label{eq: lower obund for abel proof}
% % \Gamma(k,k') d_{k} + \big(1-\Gamma(k,k')\big) d_\pi \ge \delta.
% % \end{align}
% % To bound the stochastic part, we note that for a fixed terminal index $k'$, 
% Since the map $i \mapsto \Gamma(i+1,k')$ is non-decreasing with values in $(0,1]$, using summation by parts, the triangle inequality and equation \eqref{eq: definition lock in event single gap} yield
% \begin{align}
% \nonumber
% \left|
% \sum_{i=k}^{k'-1} \Gamma(i+1,k') \alpha_i \xi_i
% \right|
% &\le
% \big| \Gamma(k',k') \big| \big| \Xi(k,k') \big| + \sum_{i=k}^{k'-2} \big| \Gamma(i+2,k') - \Gamma(i+1,k') \big| \big| \Xi(k, i+1) \big|
% \\
% \nonumber
% &\le
% \left( \Gamma(k',k') + \sum_{i=k}^{k'-2} \big( \Gamma(i+2,k') - \Gamma(i+1,k') \big) \! \right)
% \sup_{k \le m \le k'} \bigl| \Xi(k,m\mid s,a) \bigr|
% % \\
% % \nonumber
% % &=
% % \big( \Gamma(k',k') + \Gamma(k',k') - \Gamma(k+1,k') \big)
% % \sup_{k \le m \le k'} \bigl| \Xi(k,m\mid s,a) \bigr|
% \\
% \nonumber
% &\le
% 2 \sup_{k \le m \le k'} \bigl| \Xi(k,m\mid s,a) \bigr|
% \\
% \label{eq:abel equation on noise}
% &\le
% 2 \left( \bigl|\Xi(t_n,k \mid s,a)\bigr|
% +
% \sup_{k \le m \le k'} \bigl|\Xi(t_n, m \mid s,a)\bigr| \right)
% <
% \frac{\delta}{2}.
% \end{align}
% % On $\mathcal E^l_{t_n}(s,a)$, the triangle inequality also implies that for every $m > k$,
% % \begin{align}
% % \label{thm: triangle on Xi}
% % \bigl|\Xi(k,m\mid s,a)\bigr|
% % % &=
% % % \bigl|
% % % \Xi(t_n , m \mid s,a) - \Xi(t_n, k \mid s,a)
% % % \bigr|
% % \le
% % \bigl|\Xi(t_n,k\mid s,a)\bigr|
% % +
% % \bigl|\Xi(t_n,m\mid s,a)\bigr|
% % <
% % \frac{\delta}{4}.
% % \end{align}
% Hence, \eqref{eq: induction for abel proof} and
% \eqref{eq:abel equation on noise}
% show that $d_{k'} \ge \delta - \delta/2 > 0$ for every $k' \in \{k,\dots,t_{n+1}\}$,
% which proves \eqref{eq: gap remains strictly positive}.
% % This strictly limits the noise term in \eqref{eq: induction for abel proof} by
% % \begin{equation}
% % \label{eq:shifted-noise-bound}
% % \left|
% % \sum_{i=k}^{k'-1} \Gamma(i+1,k') \alpha_i \xi_{i}
% % \right|
% % < \frac{\delta}{2}.
% % \end{equation}
% % Plugging \eqref{eq: lower obund for abel proof} and \eqref{eq:shifted-noise-bound} into \eqref{eq: induction for abel proof}  
% % yields
% % \[
% % d_{k'}
% % >
% % \delta - \frac{\delta}{2}
% % >
% % 0
% % \qquad
% % \forall \, k' \in \{k,\dots,t_{n+1}\} .
% % \]
% % for every $k' \in \{k,\dots,t_{n+1}\}$, which proves \eqref{eq: gap remains strictly positive}.


% Finally, 
% % we lower-bound the probability of $\mathcal E^l_{t_n}(s,a)$.
% % We observe that 
% the complement event $\overline{\mathcal E^l_{t_n}(s,a)}$
% % implies that for some time $k \ge t_{n}$,
% % \[
% % \bigl| \Xi(t_n, k \mid s,a)\bigr|
% % \ge \frac{\delta}{8} ,
% % \]
% % or equivalently
% satisfies
% \[
% % \left\{\sup_{k\ge t_n}|\Xi(t_n,k\mid s,a)|\ge \frac{\delta}{8}\right\}
% \overline{\mathcal E^l_{t_n}(s,a)}
% \equiv
% \bigcup_{t = t_n}^\infty
% \left\{
% \max_{t_n\le k\le t}|\Xi(t_n,k\mid s,a)|
% \ge \frac{\delta}{8}
% \right\} \! .
% \]
% If $t_n \ge K_\xi$,
% % \autoref{thm: properties noise gap} indicates that there exists an almost surely finite time $K_\xi$ such that
% the sequence $\big\{\Xi(t_n, k \mid s,a)\big\}_{t_n\le k \le t}$ is a square-integrable martingale adapted to the filtration $\{\mathcal F_{k}\}_{t_n \le k\le t}$.
% Hence, Doob's martingale inequality and the monotone convergence of conditional probabilities yield
% \begin{align}
% \nonumber
% \mathbb P\!\left(
% \overline{\mathcal E^l_{t_n}(s,a)}
% \,\middle|\,
% \mathcal F_{t_n},\pi_{t_n}
% \right)
% &=
% \lim_{t \to \infty}
% \mathbb P\!\left(
% \max_{t_n \le k \le t}| \Xi(t_n, k\mid s,a) |
% \ge \frac{\delta}{8}
% \,\middle|\,
% \mathcal F_{t_n},\pi_{t_n}
% \right) \! .
% % \\
% % \label{eq:hold-doob-trunc}
% % &\le
% % \frac{64}{\delta^2}
% % \limsup_{m\to\infty}
% % \mathbb E\!\left[M_m^2\mid \mathcal F_{t_n},\pi_{t_n}\right].
% \\
% \label{eq: prob of lock in in proof}
% &\le
% \lim_{t \to \infty}
% \frac{64}{\delta^2}
% \mathbb E\!\left[\Xi(t_n,t\mid s,a)^2 \mid \mathcal F_{t_n}, \pi_{t_n} \right] \! .
% \end{align}
% % So if $t_n \ge K_\xi$,
% % % \autoref{thm: properties noise gap} indicates that there exists an almost surely finite time $K_\xi$ such that
% % the sequence $\big\{\Xi(t_n, k \mid s,a)\big\}_{t_n\le k \le t}$ is a square-integrable martingale adapted to the filtration $\{\mathcal F_{k}\}_{t_n \le k\le t}$.
% % Hence, Doob's martingale inequality yields
% % \begin{align}
% % \label{eq: prob of lock in in proof}
% % \mathbb P\!\left(
% % \max_{t_n \le k \le t} \bigl| \Xi(t_n, k \mid s,a)\bigr|
% % \ge \frac{\delta}{8}
% % \,\middle|\,
% % \mathcal F_{t_n}, \pi_{t_n}
% % \right)
% % \le
% % \frac{64}{\delta^2}
% % \mathbb E\!\left[\Xi(t_n,t\mid s,a)^2\mid \mathcal F_{t_n}, \pi_{t_n} \right] \! .
% % \end{align}
% % Using the orthogonality of martingale increments and the variance bound from \eqref{eq:xi-var}, we obtain
% Using the tower property and \autoref{thm: properties noise gap}, we obtain
% \begin{align*}
% \mathbb E\!\left[\Xi(t_n, t \mid s,a)^2\mid \mathcal F_{t_n}, \pi_{t_n} \right]
% &=
% \sum_{i=t_n}^{t-1} \alpha_{i}^2\,
% \mathbb E\!\left[\xi_{i}^2\mid \mathcal F_{t_n}, \pi_{t_n} \right]
% \\
% &=
% \sum_{i=t_n}^{t-1} \alpha_{i}^2\,
% \mathbb E\!\left[
% \mathbb E\!\left[\xi_{i}^2\mid \mathcal F_{i}, \pi_{t_n}\right]
% \middle|\mathcal F_{t_n} , \pi_{t_n}
% \right]
% \\
% &\le
% c_\xi \sum_{i=t_n}^{t-1} \alpha_{i}^2
% % \\
% % &\le
% % c_\xi \sum_{i=t_n}^{\infty} \alpha_{i}^2 .
% \end{align*}
% Substituting this back in \eqref{eq: prob of lock in in proof} proves \eqref{eq:probability-small-remaining-noise} because
% \[
% \mathbb P\!\left(
% \overline{\mathcal E^l_{t_n}(s,a)}
% \,\middle|\,
% \mathcal F_{t_n}, \pi_{t_n}
% \right)
% \le
% \frac{64 c_\xi}{\delta^2}\sum_{i=t_n}^{\infty} \alpha_{i}^2 .
% \]
% Since $n$, $(s,a)$, and $k$ were chosen arbitrarily from their respective sets, the result holds uniformly.
% \end{proof}

% % In other words, for sufficiently large $k$, whenever a gap $d^b_k$ reaches a fixed $\delta > 0$, it remains positive with large probability. 


% %%%%%%%%%%%%%%%%%%%%%%%%%%%%%%%%%%%%%%%%%%%%%%%%%%%%%%%%%%%%%%%
% \paragraph{Global event for all gaps to remain strictly positive}

% % so that none of them become greedy up to and including the transition time $t_{n+1}$.
% Fix an $n \in \mathbb Z_{\ge 0}$,
% let $L := |\mathcal B_{\pi_{t_n}}|$,
% and assume that $L>0$, otherwise \autoref{thm: fixed prob that bad actions dont become greedy} is immediate.
% % We now apply the three sequential events to every pair in $\mathcal B_{\pi_{t_n}}$.
% Conditioned on $\{\mathcal F_{t_n}, \pi_{t_n}\}$, enumerate
% \[
% \mathcal B_{\pi_{t_n}} = \big\{ (s_1,a_1), \dots, (s_L,a_L) \big\} .
% \]
% We first apply the $M$-step seed block of \autoref{thm: event to reach a_k margin} to $(s_\ell,a_\ell)$ at time $b_\ell := t_n + (\ell-1)M$ for each $\ell \in \{1,\dots,L\}$.
% At the end of this seed phase, at time $T_n := t_n + LM$, we apply the growth event of \autoref{thm: from a_k margin to delta} simultaneously to all pairs. The lock-in event of \autoref{thm: event noise remains small} is enforced throughout.
% This yields the global events
% \begin{align}
% \label{eq: def global events}
% \mathcal E^s_{t_n} &\coloneqq \bigcap_{\ell=1}^{L} \mathcal E^s_{b_\ell}(s_\ell, a_\ell),
% \qquad
% \mathcal E^g_{t_n} \coloneqq \bigcap_{\ell=1}^{L} \mathcal E^g_{T_n}(s_\ell, a_\ell),
% \qquad
% \mathcal E^l_{t_n} \coloneqq \bigcap_{\ell=1}^{L} \mathcal E^l_{t_n}(s_\ell, a_\ell).
% \end{align}

% % We first apply the $M$-step seed block from \autoref{thm: event to reach a_k margin} to pair $(s_\ell,a_\ell)$
% % at time $b_\ell := t_n + (\ell-1) M$ for $\ell \in \{1, \dots, L\}$.

% % At the end of the seed phase, i.e. at time $T_n := t_n + L M$,
% % we apply the $N(T_n)$-step growth block of \autoref{thm: from a_k margin to delta} to each pair.

% % Finally, the global lock-in event of \autoref{thm: event noise remains small} ensures

% % % to keep them strictly positive until time $t_{n+1}$.

% % Altogether, we obtain the global seed, growth and lock-in events
% % \begin{align}
% % \label{eq: def global events}
% % % \label{eq: def global seed event}
% % \mathcal E^s_{t_n} &\coloneqq \bigcap_{\ell=1}^{L} \mathcal E^s_{b_\ell}(s_\ell, a_\ell) ,
% % % \\
% % \qquad
% % % \label{eq: def global grow event}
% % \mathcal E^g_{t_n} &\coloneqq \bigcap_{\ell=1}^{L} \mathcal E^g_{T_n}(s_\ell, a_\ell) ,
% % % \\
% % \qquad
% % % \label{eq: def global hold event}
% % \mathcal E^l_{t_n} &\coloneqq \bigcap_{\ell=1}^{L} \mathcal E^l_{t_n}(s_\ell, a_\ell) .
% % \end{align}


% The gap of a pair $(s_\ell, a_\ell)$ may remain zero while $k \le b_\ell$.
% During this period, the policy $\pi_k$ may be resampled among greedy actions, and the tie can be broken in favour of $(s_\ell, a_\ell)$.
% To exclude this behaviour, we introduce a global tie-breaking event that prevents any ``bad'' action from being selected while its gap is zero:
% \begin{align}
% \label{eq: def global tie event}
% \mathcal E^t_{t_n} \coloneqq
% \bigcap_{k=t_n}^{\min \{ T_n-1 , t_{n+1} \}}
% \left\{
% \forall \, s \in \mathcal S : 
% (s, \pi_k(s)) \notin \mathcal B_{\pi_{t_n}}
% \right\} .
% \end{align}

% % We now show that, on $\mathcal E^t_{t_n} \cap \mathcal E^s_{t_n} \cap \mathcal E^g_{t_n} \cap \mathcal E^l_{t_n}$
% % no pair in $\mathcal B_{\pi_{t_n}}$ is selected up to and including time $t_{n+1}$.
% % the greedy policy at time $t_{n+1}$ does not correspond to any pair in $\mathcal B_{\pi_{t_n}}$.


% \begin{lemma}
% \label{thm: event that yields positive gaps}
% Consider the setting of \autoref{thm: fixed prob that bad actions dont become greedy} and let $K_N$ be the finite time obtained in \autoref{thm: growth window}.
% For every index $n \in \mathbb Z_{\ge 0}$ such that $t_n \ge K_N$,
% the joint event $\mathcal E^t_{t_n} \cap \mathcal E^s_{t_n} \cap \mathcal E^g_{t_n} \cap \mathcal E^l_{t_n}$ guarantees that
% % the greedy policy $\pi_k$ never coincides with a bad action, 
% % meaning
% % for every time $k \in \{t_n, \dots, t_{n+1}\}$
% % and every state $s \in \mathcal S$,
% \begin{align}
% \label{eq: proof combined event}
% (s,\pi_{k}(s)) \notin \mathcal B_{\pi_{t_n}}
% \qquad
% \forall \, s \in \sset, \
% \forall \, k \in \{t_n, \dots, t_{n+1}\} .
% \end{align}
% \end{lemma}

% \begin{proof}
% Fix an $n \in \mathbb Z_{\ge 0}$ such that $t_n \ge K_N$, and let $\pi := \pi_{t_n}$.
% For each pair $(s,a) \in \mathcal B_\pi$, define the sequence $\big\{ d_k(s,a) \big\}_{k=t_n}^{t_{n+1}}$ as in \autoref{def: gap variables}.
% % We track the evolution of a specific gap $d_k(s_\ell, a_\ell)$ across three phases.
% We show that, on the joint event, no pair in $\mathcal B_{\pi}$ is selected up to and including time $t_{n+1}$.


% Fix a pair $(s_\ell,a_\ell) \in \mathcal B_\pi$.
% At time $t_n$, all gaps are non-negative as $\pi_{t_n}$ is greedy
% with respect to $q_{t_n}$.
% Before the seed block of $(s_\ell,a_\ell)$ starts, earlier seed blocks
% cannot decrease its gap.
% Indeed, if the seed block of a pair $(s_j,a_j)$ with $j<\ell$ updates
% the pair $(s_j,a)$, then the case distinction in
% \eqref{eq: def good prefix} gives
% $(s_j,a) \neq (s_\ell,a_\ell)$.
% Thus, if $s_j \neq s_\ell$, the gap $d_k(s_\ell,a_\ell)$ is unchanged;
% if $s_j = s_\ell$ and $a \neq \pi(s_\ell)$, it is also unchanged; 
% and if $s_j=s_\ell$ and $a=\pi(s_\ell)$, it increases, as in
% \eqref{eq: final lower bound case 1} and \eqref{eq: increase in d^b case B}.
% Hence, by time $b_\ell$, either the gap has already reached $\delta$,
% or
% \[
% 0 \le d_{b_\ell}(s_\ell,a_\ell) < \delta .
% \]


% % By definition, all gaps are non-negative at time $t_n$ since $\pi_{t_n}$ is greedy with respect to $q_{t_n}$.
% % Moreover, on the joint event \(\mathcal E^t_{t_n} \cap \mathcal E^s_{t_n}\), the gap \(d_k(s_\ell,a_\ell)\) cannot decrease due to an earlier seed block \(\mathcal E^s_{b_j}(s_j,a_j)\) acting on a pair \(j<\ell\).
% % More precisely, let $(s_j, a)$ be the state-action pair updated by $\mathcal E^s_{b_j}(s_j,a_j)$. Since $j \neq \ell$ and $a_\ell \neq \pi(s_\ell)$, the case choice in \eqref{eq: def good prefix} implies that $(s_j,a) \neq (s_\ell,a_\ell)$.
% % So if $s_j \neq s_\ell$, then \(d_k(s_\ell,a_\ell)\) is unchanged;
% % if $s_j = s_\ell$ but $a \neq \pi(s_\ell)$, then \(d_k(s_\ell,a_\ell)\) is unchanged;
% % and if $s_j = s_\ell$ and $a = \pi(s_\ell)$, then \(d_k(s_\ell,a_\ell)\) increases exactly as in \eqref{eq: final lower bound case 1} and \eqref{eq: increase in d^b case B}.
% % Thus, by time \(b_\ell\),
% % either the gap associated with $(s_\ell,a_\ell)$ already hit \(\delta\), in which case \(\mathcal E^l_{t_n}(s_\ell,a_\ell)\) keeps it strictly positive up to time $t_{n+1}$ (\autoref{thm: event noise remains small}), or else \(0 \le d_{b_\ell}(s_\ell, a_\ell) < \delta\).


% If the gap has already reached $\delta$, then
% $\mathcal E^l_{t_n}(s_\ell,a_\ell)$ keeps it strictly positive up to
% time $t_{n+1}$ by \autoref{thm: event noise remains small}.
% Otherwise, applying \autoref{thm: event to reach a_k margin} at time
% $b_\ell$ on $\mathcal E^s_{b_\ell}(s_\ell,a_\ell)$ yields
% \[
% d_k(s_\ell,a_\ell) > c_s A(b_\ell,k) > 0
% \qquad
% \forall \, b_\ell < k \le \min\{ \tau^s_{b_\ell}(s_\ell,a_\ell), b_\ell+M, t_{n+1} \}.
% \]
% Thus during its seed block, the gap either reaches $\delta$, or, at the end of the block,
% \[
% d_{b_\ell+M}(s_\ell,a_\ell)
% > c_s A(b_\ell,b_\ell+M)
% \ge c_s M \alpha_{b_\ell+M}
% \ge c_s M \alpha_{T_n}
% \ge c_s \rho_\alpha M \alpha_{T_n}.
% \]
% Here we used that $\{\alpha_k\}_{k\ge0}$ is non-increasing and that
% $\rho_\alpha \in (0,1]$.


% % In the latter case, \autoref{thm: event to reach a_k margin} applied at time \(b_\ell\) on \(\mathcal E^s_{b_\ell}(s_\ell,a_\ell)\) yields
% % \[
% % d_{k}(s_\ell,a_\ell) > c_s A(b_\ell, k) > 0
% % \]
% % for every $k$ that satisfies $b_\ell < k \le \big( \tau^s_{b_\ell}(s_\ell,a_\ell) \wedge (b_\ell + M) \wedge t_{n+1} \big)$.
% % If this gap does not reach \(\delta\) during its $M$-step seed block, then
% % \[
% % d_{b_\ell+M}(s_\ell,a_\ell) 
% % > c_s A(b_\ell, b_\ell+M)
% % \ge c_s M \alpha_{b_\ell+M}
% % \ge c_s M \alpha_{T_n}
% % \ge c_s \rho_\alpha M \alpha_{T_n}
% % \]
% % by the non-increasing property of $(\alpha_k)_{k \ge 0}$, and the fact that $\rho_\alpha \in (0,1]$.
% % Since later seed blocks can also not decrease this gap,
% % either it reaches $\delta$ before time $T_n$ and then \(\mathcal E^l_{t_n}(s_\ell,a_\ell)\) keeps it strictly positive up to $t_{n+1}$,
% % or else
% % \[
% % d_{T_n}(s_\ell,a_\ell)
% % > c_s \rho_\alpha M \alpha_{T_n}.
% % \]

% Later seed blocks cannot decrease this gap by the same argument as above. Therefore, at time $T_n$, either the gap has already reached $\delta$, in which case the lock-in event keeps it strictly positive up to $t_{n+1}$, or else
% \[
% d_{T_n}(s_\ell,a_\ell)
% > c_s \rho_\alpha M \alpha_{T_n}.
% \]
% In the latter case. Since $T_n \ge t_n \ge K_N$,
% applying \autoref{thm: from a_k margin to delta} on $\mathcal E^g_{T_n}(s_\ell,a_\ell)$
% % \[
% % \mathcal E^g_{T_n}(s_\ell,a_\ell)
% % \cap
% % \big\{
% % c_s \rho_\alpha M \alpha_{T_n}
% % \le d_{T_n}(s_\ell,a_\ell) < \delta
% % \big\}
% % \]
% ensures that the gap remains strictly positive until it reaches $\delta$,
% or until $t_{n+1}$ if the transition occurs first.
% If the gap reaches $\delta$ before $t_{n+1}$, the lock-in event again keeps it strictly positive up to $t_{n+1}$.

% % Suppose the gap has not yet hit \(\delta\) by time \(T_n\). 
% % Since $T_n > t_n \ge K_N$,
% % \autoref{thm: from a_k margin to delta} guarantees that on the joint event
% % \[
% % \mathcal E^g_{T_n}(s_\ell,a_\ell) \cap
% % \big\{ c_s \rho_\alpha M \alpha_{T_n} \le d_{T_n}(s_\ell,a_\ell) < \delta \big\}
% % \]
% % that gap remains strictly positive until it reaches $\delta$ by time $T_n + N(T_n)$, or at least until the next transition time $t_{n+1}$, whichever happens first.
% % Again, if the gap reaches $\delta$ before $t_{n+1}$, then $\mathcal E^l_{t_n}(s_\ell, a_\ell)$ keeps it strictly positive up to $t_{n+1}$.

% % In summary,
% % for every $k \in \{t_n, \dots, T_n-1\}$,
% % the joint event $\mathcal E^t_{t_n} \cap \mathcal E^s_{t_n} \cap \mathcal E^l_{t_n}$
% % guarantees that $d_k(s,a) \ge 0$ for every $(s,a) \in \mathcal B_{\pi_{t_n}}$ and $(s,\pi_{k}(s)) \notin \mathcal B_{\pi_{t_n}}$ when a gap is zero in $s$.
% % As a result, if the transition occurs before time $T_n$, meaning if $t_{n+1} < T_n$, then, at time $t_{n+1}$,
% % for each pair $(s,a) \in \mathcal B_{\pi_{t_n}}$, 
% % either action $a$ is not greedy with respect to $q_{t_{n+1}}$,
% % or it is tied with $\pi_{t_{n}}(s)$ and $\mathcal E^t_{t_n}$ ensures the tie is broken in favour of an action outside $\mathcal B_{\pi_{t_n}}$, for instance $\pi_{t_{n}}(s)$.
% % In case $t_{n+1} \ge T_n$,
% % the joint event $\mathcal E^g_{t_n} \cap \mathcal E^l_{t_n}$
% % ensures that all gaps remain strictly positive
% % for the remaining time steps $k \in \{T_n, \dots, t_{n+1}\}$, and hence for every pair $(s,a) \in \mathcal B_{\pi_{t_n}}$,
% % \[
% % q_k(s, \pi_k(s)) \ge q_k(s, \pi_{t_n}(s)) > q_k(s,a).
% % \]
% % So no bad action can be greedy with respect to $q_k$ after time $T_n$.
% % Overall, the joint event $\mathcal E^t_{t_n} \cap \mathcal E^s_{t_n} \cap \mathcal E^g_{t_n} \cap \mathcal E^l_{t_n}$ guarantees that $(s,\pi_{k}(s)) \notin \mathcal B_{\pi_{t_n}}$ for all $k \in \{t_n, \dots, t_{n+1}\}$.

% Consequently, for every $(s,a) \in \mathcal B_\pi$ and every
% $k \in \{t_n,\dots,t_{n+1}\}$, either the gap $d_k(s,a)$ is zero during the seed phase, in which case $\mathcal E^t_{t_n}$ prevents the policy from selecting $(s,a)$, 
% or the gap is strictly positive,
% in which case $(s,a)$ cannot be greedy because
% \[
% q_k(s,\pi_k(s))
% \ge q_k(s,\pi(s))
% > q_k(s,a) .
% \]
% % If $t_{n+1}<T_n$, this conclusion follows from
% % $\mathcal E^t_{t_n} \cap \mathcal E^s_{t_n} \cap \mathcal E^l_{t_n}$.
% % If $t_{n+1}\ge T_n$, it follows from $\mathcal E^g_{t_n} \cap \mathcal E^l_{t_n}$.
% % Therefore, combining the events yields
% % \[
% % (s,\pi_k(s)) \notin \mathcal B_{\pi_{t_n}}
% % \qquad
% % \forall\, s \in \mathcal S , \
% % \forall\, k \in \{t_n,\dots,t_{n+1}\} .
% % \]
% % This proves the claim.
% Either way, the joint event $\mathcal E^t_{t_n} \cap \mathcal E^s_{t_n} \cap \mathcal E^g_{t_n} \cap \mathcal E^l_{t_n}$ guarantees that
% \[
% (s,\pi_k(s)) \notin \mathcal B_{\pi_{t_n}}
% \qquad
% \forall\, s \in \mathcal S , \
% \forall\, k \in \{t_n,\dots,t_{n+1}\} .
% \]
% which proves \eqref{eq: proof combined event}.
% \end{proof}

 
% %%%%%%%%%%%%%%%%%%%%%%%%%%%%%%%%%%%%%%%%%%%%%%%%%%%%%%%%
% \paragraph{Proof of \autoref{thm: fixed prob that bad actions dont become greedy}}

% \begin{proof}
% Define the constants
% \begin{align}
% p_t &:= \beta^{|\mathcal S| } \in (0,1) ,
% \\
% p_s &:= \sigma_{\min} \frac{d_{\min}^2}{d_{\min}^2 + 16 c_v} \in (0,1) ,
% \\
% p_g
% &:=
% 1-|\mathcal S\times\mathcal A|\,
% \frac{8 c_\xi}{c^2\rho_\alpha^2}
% \sum_{j=0}^{\infty}\frac{2^j}{(M+2^j)^2} ,
% \\
% p_{\overline l} &:= 
% |\mathcal S\times\mathcal A|
% \frac{64 c_\xi}{\delta^2}\sum_{i=t_n}^\infty \alpha_i^2 .
% \end{align}
% Since $\lim_{M \to \infty} \sum_{j \ge 0} 2^j/(M+2^j)^2 = 0$, we can choose $M \in \mathbb N$ sufficiently large such that $p_g \ge 1/2$.
% The Robbins--Monro conditions in \eqref{eq: robbins monro conditions} also ensure that there exists a \(K_l<\infty\) such that $p_{\overline l} \le \frac14 ( p_t p_s ) ^{|\mathcal S\times\mathcal A| M}$ for all \(t_n \ge K_l\).


% Fix an index $n \in \mathbb Z_{\ge 0}$ such that $t_n \ge \max\{K_\xi, K_N, K_l\}$
% where $K_\xi$ and $K_N$ be the finite times obtained in \autoref{thm: properties noise gap} and \autoref{thm: growth window}, respectively.
% Let $L := |\mathcal B_{\pi_{t_n}}|$ and assume that $L>0$, otherwise \autoref{thm: fixed prob that bad actions dont become greedy} is immediate.


% From \autoref{thm: event that yields positive gaps} and $\mathcal E^t_{t_n} \cap \mathcal E^s_{t_n} \in \mathcal F_{T_n}$, it follows that
% % no bad action becomes greedy. 
% % In other words, when $\pi_{t_n}$ is suboptimal, we obtain at some almost surely finite time $t_{n+1}$:
% % \begin{align}
% % \label{eq: condition that policy improves}
% % q_{\pi_{t_n}} \bigl( s, \pi_{t_{n+1}}(s) \bigr) \ge q_{\pi_{t_n}} \bigl(s,\pi_{t_n}(s) \bigr)
% % \end{align}
% % for every state $s\in S$.
% % By definition of the transition times, $q_{\pi_{t_n}} \neq q_{\pi_{t_{n+1}}}$.
% % Thus \autoref{thm: strict improvement theorem} proves that policy $\pi_{t_{n+1}}$ must be better than $\pi_{t_{n}}$, meaning $\pi_{t_{n+1}} \succ \pi_{t_n}$.
% % As a result, if $\pi_{t_n} \notin \mathcal P_*$, then
% % \[
% % \{\pi_{t_{n+1}} \succ \pi_{t_n} \}
% % \ \supseteq \
% % \mathcal E^t_{t_n} \cap \mathcal E^s_{t_n} \cap \mathcal E^g_{t_n} \cap \mathcal E^l_{t_n} .
% % \]
% % So it follows that, on the event $\{\pi_{t_n} \notin \mathcal P_*\}$,
% \begin{align}
% \nonumber
% &\mathbb P \big(
% \forall \, s \in \sset, \
% \forall \, k \in \{t_n, \dots, t_{n+1}\}
% :
% (s,\pi_{k}(s)) \notin \mathcal B_{\pi_{t_n}}
% \mid \mathcal F_{t_n}, \pi_{t_n}
% \big)
% \\
% \nonumber
% &\qquad\ge
% \mathbb P ( \mathcal E^t_{t_n} \cap \mathcal E^s_{t_n} \cap \mathcal E^g_{t_n} \cap \mathcal E^l_{t_n} \mid \mathcal F_{t_n}, \pi_{t_n} )
% % \end{multline*}
% % Since $\mathcal E^t_{t_n} \cap \mathcal E^s_{t_n} \in \mathcal F_{T_n}$, we can decompose the probability as
% % \begin{align}
% % \nonumber
% % \mathbb P \bigl(\mathcal E^t_{t_n} \cap \mathcal E^s_{t_n} \cap \mathcal E^g_{t_n} \cap \mathcal E^l_{t_n} \mid \mathcal F_{t_n} , \pi_{t_n} \bigr)
% \\
% \nonumber
% &\qquad\ge
% \mathbb P \bigl(\mathcal E^t_{t_n} \cap \mathcal E^s_{t_n} \cap \mathcal E^g_{t_n} \mid \mathcal F_{t_n} , \pi_{t_n} \bigr) 
% - \mathbb  P\bigl(\overline{\mathcal E^l_{t_n}} \mid \mathcal F_{t_n} , \pi_{t_n} \bigr)
% \\
% \label{eq: partial decomposition of the joint probabilities}
% &\qquad= \mathbb E \left[ \mathbf 1_{\mathcal E^t_{t_n} \cap \mathcal E^s_{t_n}} \mathbb P(\mathcal E^g_{t_n} \mid \mathcal F_{T_n}) \mid \mathcal F_{t_n}, \pi_{t_n} \right] - \mathbb P\bigl(\overline{\mathcal E^l_{t_n}} \mid \mathcal F_{t_n} , \pi_{t_n} \bigr)
% \end{align}
% Given that $T_n > t_n \ge \max \{K_\xi, K_N\}$, we obtain from \autoref{thm: from a_k margin to delta} and our specific choice of $M$ that
% \begin{align*}
% \mathbb P(\mathcal E^g_{t_n} \mid \mathcal F_{T_n}) 
% &\ge 1 - \sum_{(s,a) \in \mathcal B_{\pi_{t_n}}} \mathbb P \bigl( \overline{\mathcal E^g_{T_n}(s,a)} \mid \mathcal F_{T_n} \bigr)
% % \\
% % &\ge 1 - |\mathcal B_{\pi_{t_n}}| \frac{8 c_\xi}{c^2 \rho_\alpha^2} \sum_{j = 0}^{\infty} \frac{2^j}{(M+ 2^j)^2}
% \\
% &\ge 1 - |\sset \times \aset| \frac{8 c_\xi}{c^2 \rho_\alpha^2} \sum_{j = 0}^{\infty} \frac{2^j}{(M+ 2^j)^2}
% = p_g \ge \frac12 .
% \end{align*}
% Further, given that $t_n \ge \max \{K_\xi, K_l\}$, we obtain from \autoref{thm: event noise remains small} that
% \begin{align*}
% \mathbb P \bigl( \overline{\mathcal E^l_{t_n}} \mid \mathcal F_{t_n} , \pi_{t_n} \bigr)
% &\le
% \sum_{(s,a) \in \mathcal B_{\pi_{t_n}}} \mathbb P \bigl(\overline{\mathcal E^l_{t_n}(s,a)} \mid \mathcal F_{t_n}, \pi_{t_n} \bigr)
% % \\
% % &\le
% % |\mathcal B_{\pi_{t_n}}|
% % \frac{64 c_\xi}{\delta^2}\sum_{i={t_n}}^\infty \alpha_i^2
% \\
% &\le
% |\mathcal S\times\mathcal A|
% \frac{64 c_\xi}{\delta^2}\sum_{i={t_n}}^\infty \alpha_i^2
% =
% p_{\overline l}
% \le
% \frac14 ( p_t p_s)^{|\mathcal S \times \mathcal A| M} .
% \end{align*}
% Thus we can further bound \eqref{eq: partial decomposition of the joint probabilities} by
% \begin{multline}
% \label{eq: probability lower bound intermediary}
% % \mathbb  P(\mathcal E^t_{t_n} \cap \mathcal E^s_{t_n} \cap \mathcal E^g_{t_n} \cap \mathcal E^l_{t_n} \mid \mathcal F_{t_n}, \pi_{t_n} )
% \mathbb P \big(
% \forall \, s \in \sset, \
% \forall \, k \in \{t_n, \dots, t_{n+1}\}
% :
% (s,\pi_{k}(s)) \notin \mathcal B_{\pi_{t_n}}
% \mid \mathcal F_{t_n}, \pi_{t_n}
% \big)
% \\
% \ge
% p_g \, \mathbb P(\mathcal E^t_{t_n} \cap \mathcal E^s_{t_n} \mid \mathcal F_{t_n}, \pi_{t_n}) - p_{\overline l} .
% \end{multline}
% We now lower-bound $\mathbb P(\mathcal E^t_{t_n} \cap \mathcal E^s_{t_n} \mid \mathcal F_{t_n}, \pi_{t_n})$ by sequentially conditioning over $\{t_n, \dots, T_n-1\}$.
% For this, define
% \[
% \ell(k) := 1 + \left\lfloor \frac{k-t_n}{M} \right\rfloor \in \{1,\dots,L\} ,
% \]
% so that \(k\) lies in the \(M\)-step seed block of pair \((s_{\ell(k)},a_{\ell(k)})\).
% % starting at time $b_{\ell(k)} = t_n + (\ell(k)-1) M$.
% We decompose events \(\mathcal E^s_{t_n}\) and \(\mathcal E^t_{t_n}\) into one-step events as follows. For $k \in \{t_n,\dots,T_n-1\}$, let
% \[
% \mathcal E^s_{t_n}(k)
% :=
% \left\{
% u_{k}(s_{\ell(k)},\pi_{t_n}(s_{\ell(k)}))=1,\ 
% v_{k}(s_{\ell(k)},\pi_{t_n}(s_{\ell(k)})) \ge -\dfrac{d_{\min}}{4}
% \right\}
% \]
% if $q_{\pi_{t_n}} \big(s_{\ell(k)},\pi_{t_n}(s_{\ell(k)}) \big) - q_{b_{\ell(k)}} \big(s_{\ell(k)},\pi_{t_n}(s_{\ell(k)}) \big) \ge d_{\min}/2$, 
% and
% \[
% \mathcal E^s_{t_n}(k)
% :=
% \left\{
% u_{k}(s_{\ell(k)},a_{\ell(k)})=1,\ 
% v_{k}(s_{\ell(k)},a_{\ell(k)}) \le \dfrac{d_{\min}}{4}
% \right\}
% \]
% otherwise.
% Similarly, define
% \[
% \mathcal E^t_{t_n}(k)
% :=
% \left\{
% \forall \, s \in \mathcal S : 
% (s,\pi_k(s)) \notin \mathcal B_{\pi_{t_n}}
% \right\}
% \]
% if $k \le t_{n+1}$, and
% \[
% \mathcal E^t_{t_n}(k)
% := \Omega
% \]
% otherwise,
% where $\Omega$ is the certain event.
% These one-step events yield the decompositions
% \begin{align*}  
% \mathcal E^s_{t_n} = \bigcap_{k=t_n}^{T_n-1}\mathcal E^s_{t_n}(k),
% \qquad
% \mathcal E^t_{t_n} = \bigcap_{k=t_n}^{T_n-1}\mathcal E^t_{t_n}(k).
% \end{align*}


% Fix a time $k \in \{t_n, \dots, T_n\}$ and consider the joint event
% \begin{align}
% \label{eq: joint intermediate event}
% \bigcap_{i=t_n}^{k-1} \bigl( \mathcal E^t_{t_n}(i)\cap \mathcal E^s_{t_n}(i) \bigr).
% \end{align}
% We first bound the tie-breaking event. If $k \le t_{n+1}$, then all gaps remain non-negative up to time $k$. A pair $(s,a) \in \mathcal B_{\pi_{t_n}}$ can thus be greedy only if its gap is zero. In that case, $\pi_{t_n}(s)$ is also greedy. So, on \eqref{eq: joint intermediate event}, $\mathcal E^t_{t_n}(k)$ holds whenever the policy $\pi_k$ selects $\pi_{t_n}(s)$ in each such state. By the tie-breaking condition in \eqref{eq:greedy-pol-tie-breaking} and the independence across states,
% \[
% \mathbb P\!\left(\mathcal E^t_{t_n}(k)\mid \mathcal F_k,\pi_{t_n}\right)
% \ge \beta^{|\mathcal S|} .
% \]
% If $k>t_{n+1}$, then $\mathcal E^t_{t_n}(k)=\Omega$, so the same bound holds. Thus,
% \[
% \mathbb P\!\left(\mathcal E^t_{t_n}(k)\mid \mathcal F_k,\pi_{t_n}\right) \ge p_t.
% \]
% % Regarding the tie-breaking event, we distinguish two cases.
% % First, suppose that $k \le t_{n+1}$. 
% % Every gap associated with a pair in \(\mathcal B_{\pi_{t_n}}\) is non-negative at time \(k\) because all gaps are non-negative at time \(t_n\) and seed blocks cannot decrease any gap below zero.
% % Therefore, if a pair \((s,a)\in\mathcal B_{\pi_{t_n}}\) is greedy at time \(k\), then necessarily
% % \[
% % d_k(s,a)=q_k(s,\pi_{t_n}(s))-q_k(s,a)=0.
% % \]
% % Hence, \(\pi_{t_n}(s)\) is also greedy in state \(s\). It follows that, on the joint event \eqref{eq: joint intermediate event}, the tie-breaking condition \(\mathcal E^t_{t_n}(k)\) is satisfied whenever the policy at time \(k\) follows action \(\pi_{t_n}(s)\) in every state \(s\) where a bad action is greedy. By the action-selection condition in \eqref{eq:greedy-pol-tie-breaking} and the conditional independence of tie-breaking across states, we obtain
% % \[
% % \mathbb P\!\left(\mathcal E^t_{t_n}(k)\mid \mathcal F_k,\pi_{t_n}\right)
% % \ge
% % \beta^{|\mathcal S|}
% % =
% % p_t .
% % \]
% % Second, if \(k>t_{n+1}\), then $\mathcal E^t_{t_n}(k)$ is the certain event. So we also have
% % \[
% % \mathbb P\!\left(\mathcal E^t_{t_n}(k)\mid \mathcal F_k,\pi_{t_n}\right)=1\ge p_t .
% % \]
% % % In both cases,
% % % \[
% % % \mathbb P\!\left(\mathcal E^t_{t_n}(k)\mid \mathcal F_k,\pi_{t_n}\right)\ge p_t .
% % % \]
% For the seed event, the exploring starts condition together with
% \eqref{eq:noise-lower}--\eqref{eq:noise-upper} yields
% \[
% \mathbb P\!\left(\mathcal E^s_{t_n}(k)\mid \mathcal F_k, \pi_{t_n}\right)
% \ge \sigma_{\min}\frac{d_{\min}^2}{d_{\min}^2+16c_v}
% = p_s.
% \]
% % Regarding the seed event, the strict exploring starts conditions and inequalities \eqref{eq:noise-lower} and \eqref{eq:noise-upper} yield
% % \[
% % \mathbb P\!\left(\mathcal E^s_{t_n}(k)\mid \mathcal F_{k}, \pi_k, \pi_{t_n} \right)
% % \ge
% % \sigma_{\min}\frac{d_{\min}^2}{d_{\min}^2+16c_v}
% % =
% % p_s .
% % \]
% Combining both bounds and using the tower property, we obtain
% \begin{align*}
% \mathbb P\!\left(\mathcal E^t_{t_n}(k) \cap \mathcal E^s_{t_n}(k)\mid \mathcal F_k, \pi_{t_n}\right)
% \ge
% p_s\,\mathbb P\!\left(\mathcal E^t_{t_n}(k) \mid \mathcal F_k,\pi_{t_n}\right)
% \ge p_t p_s.
% \end{align*}
% Applying this recursively yields
% \begin{align*}
% &\mathbb P \! \left(
% \bigcap_{i=t_n}^{k} \bigl( \mathcal E^t_{t_n}(i) \cap \mathcal E^s_{t_n}(i) \bigr)
% \ \middle| \ \mathcal F_{t_n}, \pi_{t_n}
% \right)
% % \\
% % &\qquad\ge
% \ge
% p_t p_s\,
% \mathbb P\!\left(
% \bigcap_{i=t_n}^{k-1} \bigl( \mathcal E^t_{t_n}(i)\cap \mathcal E^s_{t_n}(i) \bigr)
% \ \middle|\ \mathcal F_{t_n}, \pi_{t_n}
% \right) \!.
% \end{align*}
% Iterating from $t_n$ to $T_n-1=t_n+LM-1$ gives
% \[
% \mathbb P(\mathcal E^t_{t_n} \cap \mathcal E^s_{t_n} \mid \mathcal F_{t_n}, \pi_{t_n})
% \ge (p_t p_s)^{LM}
% \ge (p_t p_s)^{|\mathcal S\times\mathcal A|\,M}.
% \]


% % Thus by using the tower property, we obtain
% % \begin{align*}
% % &\mathbb P\!\left(\mathcal E^t_{t_n}(k)\cap \mathcal E^s_{t_n}(k)\mid \mathcal F_{k}, \pi_{t_n }\right)
% % \\
% % &\qquad =
% % \mathbb E\!\left[ \mathbb E\!\left[ \mathbf 1_{\mathcal E^t_{t_n}(k)} \mathbf 1_{\mathcal E^s_{t_n}(k)} \mid \mathcal F_{k}, \pi_k, \pi_{t_n} \right] \middle| \mathcal F_k, \pi_{t_n}\right]
% % \\
% % &\qquad =
% % \mathbb E\!\left[ \mathbf 1_{\mathcal E^t_{t_n}(k)} \mathbb P \big( \mathcal E^s_{t_n}(k) \mid \mathcal F_{k}, \pi_k, \pi_{t_n} \big) \middle| \mathcal F_k, \pi_{t_n}\right]
% % \\
% % &\qquad \ge
% % p_s
% % \mathbb P\!\left(\mathcal E^t_{t_n}(k) \mid \mathcal F_{k}, \pi_{t_n} \right)
% % \ge
% % p_t p_s
% % \end{align*}
% % on $\bigcap_{i=t_n}^{k-1} \big( \mathcal E^t_{t_n}(i) \cap \mathcal E^s_{t_n}(i) \big)$.
% % Repeated conditioning in chronological order then yields
% % \begin{align*}
% % &\mathbb P\!\left(
% % \bigcap_{i=t_n}^{k} \bigl( \mathcal E^t_{t_n}(i) \cap \mathcal E^s_{t_n}(i) \bigr)
% % \ \middle|\ \mathcal F_{t_n}, \pi_{t_n}
% % \right)
% % \\
% % &\qquad=
% % \mathbb E\!\left[
% % \mathbf 1_{\cap_{i=t_n}^{k-1}(\mathcal E^t_{t_n}(i) \cap \mathcal E^s_{t_n}(i))}
% % \,
% % \mathbb P\!\left(\mathcal E^t_{t_n}(k)\cap \mathcal E^s_{t_n}(k) \mid \mathcal F_{k}, \pi_{t_n} \right)
% % \ \middle|\ \mathcal F_{t_n}, \pi_{t_n}
% % \right]
% % \\
% % &\qquad\ge
% % p_t p_s\,
% % \mathbb P\!\left(
% % \bigcap_{i=t_n}^{k-1} \bigl( \mathcal E^t_{t_n}(i) \cap \mathcal E^s_{t_n}(i) \bigr)
% % \ \middle|\ \mathcal F_{t_n}, \pi_{t_n}
% % \right).
% % \end{align*}
% % By iterating this inequality from time $t_n$ up to $T_n-1 = t_n + L M -1$,
% % and using $0 < p_t p_s \le 1$ along with $L \le |\mathcal S \times \aset|$,
% % we obtain 
% % \[
% % \mathbb P(\mathcal E^t_{t_n} \cap \mathcal E^s_{t_n} \mid \mathcal F_{t_n}, \pi_{t_n} )
% % \ge
% % (p_t p_s)^{LM}
% % \ge
% % (p_t p_s)^{|\mathcal S\times\mathcal A|\,M} .
% % \]

% % Plugging this into \eqref{eq: probability lower bound intermediary},
% Altogether, define
% \begin{align*}
% % \label{eq: def lower bound probability} 
% p
% :=
% \frac14
% (p_t \, p_s)^{|\mathcal S \times \mathcal A| \, M}
% % \\
% % \nonumber
% =
% \frac14
% \left(
% \beta^{|\mathcal S|}
% \sigma_{\min}\frac{d_{\min}^2}{d_{\min}^2+16 c_v}
% \right)^{|\mathcal S\times\mathcal A| \, M}
% >0  .
% \end{align*}
% Then for every \mcopi{} solution there exists an almost surely finite time $K \ge \max\{K_\xi, K_N, K_l\}$ such that for every transition time $t_n$ satisfying $K \le t_n < \infty$,
% \begin{align*}
% \mathbb P \big(
% \forall \, s \in \sset, \
% \forall \, k \in \{t_n, \dots, t_{n+1}\}
% :
% (s,\pi_{k}(s)) \notin \mathcal B_{\pi_{t_n}}
% \mid \mathcal F_{t_n}, \pi_{t_n}
% \big)
% \ge
% p_g (p_t p_s)^{|\mathcal S\times\mathcal A|\,M} - p_{\overline l}
% \ge
% p .
% \end{align*}
% Since this bound is uniform in $\pi_{t_n}$, \eqref{eq: porb better action value} follows.
% \end{proof}
% % \begin{align*}
% % \mathbb P\!\left(
% % \forall \, s \in \sset, \
% % \forall \, k \in \{t_n, \dots, t_{n+1}\}
% % :
% % (s,\pi_{k}(s)) \notin \mathcal B_{\pi_{t_n}}
% % \mid \mathcal F_{t_n}
% % \right)
% % &=
% % \mathbb E \! \left[ 
% % \mathbb P\!\left(
% % \forall \, s \in \sset, \
% % \forall \, k \in \{t_n, \dots, t_{n+1}\}
% % :
% % (s,\pi_{k}(s)) \notin \mathcal B_{\pi_{t_n}}
% % \mid \mathcal F_{t_n}, \pi_{t_n}
% % \right)
% % \middle| \mathcal F_{t_n}, \pi_{t_n} \notin \mathcal P_* \right]
% % \\
% % &\ge
% % \mathbb E [ p \mid \mathcal F_{t_n}, \pi_{t_n} \notin \mathcal P_* ]
% % \\
% % &= p .
% % \end{align*}
% % Since this lower bound is uniform over all realisations of $\pi_{t_n}$, we obtain \eqref{eq: porb better action value}.


% % \autoref{thm: fixed prob that bad actions dont become greedy} established that, under uniform updates over all actions in each state, the probability of a policy improvement is eventually bounded below by a fixed constant. Since this lower bound is uniform over all realisations of the greedy policy, it remains valid when the current policy is optimal. In that case, however, no improvement is possible. Instead, the same joint event guarantees that no further transition occurs.

% % \begin{proposition}
% % \label{thm: prob policy remains optimal}
% % Consider the setting of \autoref{thm: fixed prob that bad actions dont become greedy}.
% % The constant $p>0$ and random time $K<\infty$ of \autoref{thm: fixed prob that bad actions dont become greedy} ensure that
% % for every index $n \in \mathbb Z_{\ge 0}$ such that $t_n \ge K$,
% % \[
% % \mathbb P (t_{n+1} = \infty \mid \mathcal F_{t_n}, \pi_{t_n} \in \mathcal P_* ) \ge p .
% % \]
% % \end{proposition}
% % \begin{proof}
% % Fix an $n \in \mathbb Z_{\ge 0}$ such that $t_n \ge K$,
% % and assume that $\pi_{t_n}$ is optimal.
% % By definition of optimality, 
% % $q_{\pi_{t_n}}(s, \pi_{t_n}(s) ) = \max_a q_{\pi_{t_n}}(s,a)$
% % for every state $s \in \mathcal S$.
% % As a result, for every pair
% % $(s,a) \notin \mathcal B_{\pi_{t_n}}$,
% % \begin{align}
% % \label{eq: implication of reached otpimal policy}
% % q_{\pi_{t_n}}(s,a) = q_{\pi_{t_n}}(s, \pi_{t_n}(s)) .
% % \end{align}
% % Now suppose for contradiction that $t_{n+1} < \infty$
% % on the joint event $\mathcal E^t_{t_n} \cap \mathcal E^s_{t_n} \cap \mathcal E^g_{t_n} \cap \mathcal E^h_{t_n}$, as defined in \eqref{eq: def global seed event}--\eqref{eq: def global tie event}.
% % \autoref{thm: event that yields positive gaps} establishes that $(s,\pi_{t_{n+1}}(s)) \notin \mathcal B_{\pi_{t_n}}$ for every $s \in \mathcal S$. Hence, as in \eqref{eq: implication of reached otpimal policy},
% % \[
% % q_{\pi_{t_n}}(s,\pi_{t_{n+1}}(s)) = q_{\pi_{t_n}}(s,\pi_{t_n}(s)) .
% % \]
% % By \autoref{def: transition times}, $q_{\pi_{t_{n+1}}} \neq q_{\pi_{t_n}}$. Thus \autoref{thm: strict improvement theorem} implies that policy $\pi_{t_{n+1}}$ is strictly better than $\pi_{t_{n}}$, which was already optimal. Given this is impossible, $t_{n+1}$ must be almost surely infinite on the joint event.
% % Consequently, as in the proof of \autoref{thm: fixed prob that bad actions dont become greedy},
% % \[
% % \mathbb P (t_{n+1} = \infty \mid \mathcal F_{t_n}, \pi_{t_n} \in \mathcal P_* )
% % \ge
% % \mathbb P(\mathcal E^t_{t_n} \cap \mathcal E^s_{t_n} \cap \mathcal E^g_{t_n} \cap \mathcal E^h_{t_n} \mid \mathcal F_{t_n}, \pi_{t_n} \in \mathcal P_* )
% % \ge p
% % \]
% % with $p$ defined in \eqref{eq: def lower bound probability}.
% % \end{proof}




% % \[
% % \mathbb P\!\left(
% % \bigcap_{i=t_n}^{T_n-1 \wedge t_{n+1}-1} \bigl( \mathcal E^t_{t_n}(i) \cap \mathcal E^s_{t_n}(i) \bigr)
% % \ \middle|\ \mathcal F_{t_n}, \pi_{t_n}
% % \right)
% % \]
% % If the transition occurs any time before $T_n$, meaning $t_{n+1} < T_n$, then $\mathcal E^t_{t_n}(t_{n+1})$ guarantees that $(s, \pi_{t_{n+1}}(s)) \notin \mathcal B_{\pi_{t_n}}$ for every state $s \in \mathcal S$.


% % Altogether, the proofs of \autoref{thm: fixed prob that bad actions dont become greedy} and \autoref{thm: prob policy remains optimal}
% % rely on three successive events.
% % These events ensure that stochastic perturbations cannot consistently
% % undermine the policy improvement driven by the mean-field dynamics
% % by forcing ``bad'' actions to become greedy.

% % % use advantage instead of gap?

% % The lock-in phenomenon is a direct application of \citep[Theorem~2.1]{Borkar2002OnApproximation}:
% % if the gap variable for a ``bad'' pair reaches a fixed positive margin infinitely often, 
% % it must eventually remain strictly positive.
% % However, since decreasing learning rates cause these gaps to be increasingly smaller at initialization,
% % we introduced a growth event to ensure they still reach this fixed margin.
% % This event extends the lock-in phenomenon to hold under a milder condition: gaps only need to reach a decreasing margin of size $O(\alpha_k)$ infinitely often.
% % Still, it requires an additional comparability condition on the learning rates so that, with a fixed probability, once a gap hits the $O(\alpha_k)$ margin, the remaining learning rates are sufficiently large to bring it up to the original fixed margin.
% % Finally, to guarantee gaps reach the decreasing $O(\alpha_k)$ margin within a fixed horizon, we used a seed event.
% % The length of this horizon determines the probability of the growth event.
% % By fixing a sufficiently long horizon, we can drive the probability of the growth event arbitrarily close to one, regardless of the unknown parameters in \eqref{eq: prob from a_k reach delta}.
% % The only trade-off for this extended horizon is a reduced probability of the seed event.
% % Yet, since that probability is fixed once the horizon is fixed, the overall probability of improvement remains fixed.


% % It is crucial for the seed event that updates use the initial-visit method to guarantee that all gaps initially increase.
% % % only one action be updated in any state at every step so 
% % % and that potential correlation in the episode noise does not break \eqref{eq: increase in d^b case B}.
% % % The structure of the stochastic perturbations $w_k$ is crucial for the seed event.
% % % Since the stochastic perturbations caused by episode simulations can be adversarially correlated, meaning they could 
% % % (which is not really an issue because uniform updates require initial-visit anyway in practice)
% % % to exactly cancel out the update in all but one state-action pair at every step (initial-visit).
% % However, uniform update distributions are irrelevant for this part of the proof because the strict exploring starts guarantee that we can always build the seed event.
% % On the contrary, the growth and lock-in events essentially enforce that the MC-O-PI solution behaves well because it does not wander away too much from the mean-field solution.
% % Thus they rely strongly on the uniformity of the update distributions that ensures the mean-field solution behaves well in the first place.
% % The update function is, however, irrelevant.
% % All that is necessary is that the stochastic perturbation $w_k$ is unbiased and that its variance is eventually bounded.


% %%%%%%%%%%%%%%%%%%%%%%%%%%%%%%%%%%%%%%%%%%%%%%%%%%%%%%
% %%%%%%%%%%%%%%%%%%%%%%%%%%%%%%%%%%%%%%%%%%%%%%%%%%%%%%
% \subsection{Proof of \autoref{thm: iv mcopi with 1/g scaling}}

% % \begin{proposition}
% % \label{thm: convergence mcopi uniform}
% %     Solutions of \mcopi{}, as defined in \eqref{eq: mcopi sa},
% %     converge almost surely to the optimal action values $\qopt$ under \autoref{asp: standing},
% %     if
% %     the learning rates satisfy the comparability assumption (\autoref{asp: additional conditions on learning rate}),
% %     the update functions are initial-visit,
% %     and the update distributions are uniform over all actions in each state (\autoref{asp: unbiased over policies}).
% % \end{proposition}


% \begin{proof}
% Let $\{q_k, \pi_k\}_{k \ge 0}$ be an \mcopi{} solution with transition times $\{t_n\}_{n \ge 0}$.
% By \autoref{thm: fixed prob that bad actions dont become greedy}, there exists an almost surely finite time $K$ such that, with a fixed probability $p$,
% \begin{align}
% \label{eq: event no bad gaps become greedy}
% (s,\pi_k(s)) \notin \mathcal B_{\pi_{t_n}}
% \qquad
% \forall \, K \le t_n < \infty , \
% \forall \, t_n \le k \le t_{n+1} ,
% \forall \, s \in \mathcal S . \
% \end{align}
% On this event, if $t_{n+1}<\infty$, then
% \begin{align}
% \label{eq: goal of proof stoch improvement}
% q_{\pi_{t_n}} ( s, \pi_{t_{n+1}}(s) )
% \ge
% q_{\pi_{t_n}} ( s,\pi_{t_n}(s) )
% \qquad \forall\, s \in \mathcal S ,
% \end{align}
% so $\pi_{t_{n+1}}$ improves upon $\pi_{t_n}$ by \autoref{thm: strict improvement theorem}.
% If instead $t_{n+1}=\infty$, then $\pi_{t_n}=\pi_*$ by \autoref{thm: suboptimal blocks finite}.
% % and $q_k \to q_{\pi_*}$ by \cite[Proposition~4.1, Example~4.3]{Bertsekas1996Neuro-DynamicProgramming}.
% Hence, from any $t_n \ge K$, the policy either improves or is already optimal and remains so with probability at least $p$.

% Since a policy sequence can improve at most $|\mathcal P|-1$ times before reaching optimality, 
% the probability of consecutive improvements until the policy becomes optimal and then remains so is at least $p^{|\mathcal P|}$.
% Equivalently, for every $n$ with $K\le t_n<\infty$,
% \begin{align}
% \label{eq:prob-lock-within-N-transitions}
% \mathbb P\!\left(
% t_{n+|\mathcal P|} < \infty
% \mid
% \mathcal F_{t_n}
% \right)
% \le 1 - p^{|\mathcal P|}.
% \end{align}

% Fix $n \in \mathbb Z_{\ge 0}$ such that $K\le t_n<\infty$. 
% Iterating \eqref{eq:prob-lock-within-N-transitions} over blocks of length $|\mathcal P|$ gives, 
% for every $i \ge 1$,
% \begin{align*}
% \mathbb P\big(
% t_{n+i|\mathcal P|} < \infty
% \mid
% \mathcal F_{t_n}
% \big)
% &=
% \mathbb E\!\left[
% \mathbf 1_{\{t_{n+(i-1)|\mathcal P|}<\infty\}}
% \,
% \mathbb P\big(
% t_{n+i |\mathcal P|} < \infty
% \mid
% \mathcal F_{t_{n+(i-1)|\mathcal P|}}
% \big)
% \ \middle| \
% \mathcal F_{t_{n}}
% \right]
% \\
% &\le
% (1-p^{|\mathcal P|})\,
% \mathbb P (
% t_{n+(i-1) |\mathcal P|} < \infty
% \mid
% \mathcal F_{t_{n}}
% )
% \\
% &\le
% \big(1-p^{|\mathcal P|}\big)^i .
% \end{align*}
% Since $p^{|\mathcal P|} > 0$, the right-hand side converges to zero.
% Therefore,
% \[
% \mathbb{P}\!\left(
% \bigcap_{i=1}^\infty
% \big\{ t_{n+i|\mathcal P|}<\infty \big\}
% \;\middle|\;
% \mathcal F_{t_n}
% \right)
% =0.
% \]
% Thus, almost surely, there exists a finite $i_*$ such that $t_{n+i_*|\mathcal P|} = \infty$.
% As a result, there is a transition index $n_* < n + i_* |\mathcal P|$ such that $t_{n_*} < \infty$ and $t_{n_* +1} = \infty$.
% It follows that $\pi_{t_{n_*}} = \pi_*$ by \autoref{thm: suboptimal blocks finite} and $q_k \to q_{\pi_*}$ by \cite[Proposition~4.1, Example~4.3]{Bertsekas1996Neuro-DynamicProgramming}.
% \end{proof}


% \textit{ChapGPT 5.5 and Gemini 3.1 were used to structure and refine this argument.}